% English translation, made from our own transcription (original.tex), not from any existing
% translation. Every math expression, symbol, \origpage{N} marker, \tag and \emph run is
% reproduced unchanged; only the prose is translated. The memoir carries no figure, so
% \rmfigure never arises. NOTE: this header deliberately contains no dollar-delimited math,
% because pipeline/texcompare.py reads comment lines too and would count any formula quoted
% here as an extra math span.
%
% WHAT IS AND IS NOT TRANSLATED INSIDE MATH. The work's \text{...} inserts are of three kinds,
% and HOUSESTYLE R21 (prose inside a formula is translated) separates them:
%   * TRANSLATED, twice. The connective "ou" between the two alternative values of P_2 in
%     article 31 becomes \text{or}, and the parenthetical gloss on omega under display (2) of
%     article 33 becomes \text{being the surabundance}. In both the insert keeps its position
%     and every symbol around it is copied verbatim.
%   * NOT TRANSLATED: the superscript series designations of the journals cited in the
%     footnotes (French 4^e s^e, Italian 2^a, and the "t^e" of the Journal de Mathematiques).
%     These are bibliography, not prose, and stand unchanged with the rest of each citation --
%     following clebsch-1868-surfaces-algebriques, whose English translation likewise leaves its
%     German book title and the "und" joining its authors alone.
%   * NOT TRANSLATED: the authors' own cross-reference n^o to their numbered articles ("n^o 9",
%     and the plural "n^os 4, 5"). The superscript is a math span that texcompare requires to be
%     present, so "No. 9" is not available, and the printed form reads without difficulty.
%   * "etc." inside the display of article 41 is already English and is left as it stands.
%
% NAMES, HONORIFICS AND CITATIONS. Author names are letterspaced in the print and therefore
% \emph'd (R20); the \emph runs, their boundaries, and the honorific outside them are kept
% exactly as the transcription has them. "M." and "MM." are left as printed, as in the
% translations of abel-1841-fonctions-transcendantes and clebsch-1868-surfaces-algebriques.
% Max Noether's name keeps the print's spelling \emph{Nöther} (R12, notation.md). Titles of
% cited works -- French, Italian and German alike -- are bibliography and stand untranslated;
% only the words around them ("Voir les deux Mémoires", "voir aussi", "Consulter aussi sur ce
% sujet") are English here.
%
% TERMINOLOGY. Period English, and where the Italian school's English contemporaries have the
% word, theirs:
%   * genre = "genus"; genre géométrique / numérique = "geometric / numerical genus"; second and
%     third genres = "second" and "third genus"; bigenre = "bigenus".
%   * série linéaire = "linear series"; système linéaire = "linear system"; groupe (of the points
%     cut on a curve) = "group", never "divisor" -- the authors have no divisor, and the
%     anachronism would be a modernization of the notation, not of the prose.
%   * courbe adjointe, système adjoint, surface sous-adjointe = "adjoint curve", "adjoint
%     system", "sub-adjoint surface"; the authors' own Italian gloss (subaggiunta) is kept, as
%     are Nöther's German terms \emph{Restsatz}, \emph{ausgezeichnete Curven}, \emph{adjungirte
%     Flächen}, \emph{Vollschaar}, \emph{Flächengeschlecht}, \emph{Curvengeschlecht},
%     \emph{Uebergangscurve} -- these are the authors' citations of another author's words
%     (HOUSESTYLE R9).
%   * théorème du reste = "theorem of the remainder", the authors' own name for their addition
%     theorem; where they mean Brill and Nöther's theorem they write \emph{Restsatz}, and that
%     stands as printed.
%   * série caractéristique = "characteristic series"; série / système canonique = "canonical
%     series / system"; système i-canonique = "i-canonical system".
%   * défaut = "defect"; surabondance = "surabundance"; spécialité = "speciality"; dimension
%     virtuelle / effective = "virtual / effective dimension".
%   * surface régulière / irrégulière = "regular / irregular surface"; the defect that measures
%     the difference of the two genera is named by the authors' own words, not by the later
%     "irregularity".
%   * champ algébrique = "algebraic field"; champ de rationnalité = "field of rationality";
%     corps de fonctions algébriques = "field of algebraic functions" -- the three are kept
%     parallel, as the print keeps them.
%   * transformation birationnelle (ou de Cremona) = "birational (or Cremona) transformation";
%     faisceau = "pencil"; points-base = "base points"; courbe gauche = "skew curve", the period
%     English for a non-plane space curve.
%   * plan double / simple = "double / simple plane"; points tacnodals = "tacnodal points";
%     point de rebroussement = "cusp" (Salmon's word, as in the glossary of
%     clebsch-1864-anwendung-abelschen-functionen).
%   * courbe de diramation = "branch curve". Here the literal-rendering preference of R9 does
%     not apply: "diramation" is ordinary French vocabulary for branching, not a coined term of
%     art, "branch curve" is the genuine period English, and the authors' own German gloss
%     (\emph{Uebergangscurve}) stands beside it at first use.
%   * dimension / degré of a system, ordre of a curve or surface stay "dimension", "degree",
%     "order" throughout; variété algébrique = "algebraic variety"; espace ordinaire = "ordinary
%     space"; "Mémoire", capitalized in the print, is "Memoir", capital kept.
%
% PRINTER'S SLIPS. The original's inconsistent accents ("rélation", "géomètrique", "géomètrie")
% and its two French typos ("ou aura" for "on aura" in article 24, "courbe" for "courbes" in
% article 31) are kept in the transcription under R4/R23 but have no English counterpart; they
% are simply translated as the words they are, and remain visible in original.tex. The slips
% that ARE math -- the stray closing paren of "|C' + D)|" in article 15, and "3p - 6" for
% "3p_g - 6" in article 31 -- are math and are therefore copied verbatim, errors and all. The
% hyphen the print sets for a full stop after "résultent complètes" in article 27 is likewise
% carried over as printed.
%
% GUILLEMETS AND OTHER TYPOGRAPHY are the source edition's (R22) and stand where the
% transcription puts them. Em-dashed ranges (1863---1865), the bibliographic ellipsis (...) and
% the footnote series *) **) ***) †) ††) are reproduced exactly, each note remaining a complete
% unit at the end of its page's text (R15). The dateline on p316 keeps its \emph{Florence} and
% its date, set in the English order without the print's ordinal point.
\documentclass{article}
\usepackage{readmasters}

\begin{document}

\origpage{241}

\section*{On some recent results in the theory of algebraic surfaces.}

\subsection*{By MM.\ G.\ Castelnuovo in Rome and F.\ Enriques in Bologna.}

The \emph{Geometry on a general algebraic surface}, whose origin goes back to a remark of \emph{Clebsch}${}^{*)}$ and to a fundamental Memoir of M.\ \emph{Nöther}${}^{**)}$, has been enriched in these latter times with new results, thanks to the researches of the French and Italian geometers. These researches have one and the same point of departure in the cited Memoir of M.\ \emph{Nöther}, but they follow two different directions, which we shall set out briefly.

In France it is above all the transcendental functions attached to the equation of an algebraic surface that have been considered. M.\ \emph{Nöther} had already introduced into the study of surfaces certain double integrals (with algebraic differentials) which remain everywhere finite, and which play the same role as the Abelian integrals of the first kind relative to algebraic curves. Some years later M.\ \emph{Picard}${}^{***)}$ gave a new impulse to this theory by considering, alongside the double integrals, certain simple integrals of total differentials which are likewise worthy of remark. He studied in particular the simple integrals which remain everywhere finite. The best known surfaces do not truly possess such integrals; but M.\ \emph{Picard} has pointed out classes of surfaces whose remarkable properties depend precisely on the presence of transcendental functions of the kind named. M.\ \emph{Poincaré} with a short Note${}^{\dagger)}$, and M.\ \emph{Humbert}

\textbf{*)} Comptes Rendus de l'Ac.\ d.\ Sc.\ 1868.

\textbf{**)} \emph{Zur Theorie des eindeutigen Entsprechens}... I, II, Math.\ Annalen Bd.\ 2, 8, 1869, 1874.

\textbf{***)} See the two Memoirs \emph{Sur les intégrales de différentielles totales algébriques}..., Journal de Mathém.\ 4${}^{\text{e}}$ s${}^{\text{e}}$, t.\ I, 1885; the \emph{Mémoire sur la théorie des fonctions algébriques}, J.\ d.\ M.\ 4${}^{\text{e}}$ s${}^{\text{e}}$, t.\ V; as well as several Notes in the Comptes Rendus de l'Ac.\ d.\ Sc., Déc.\ 1884.

\textbf{†)} Comptes Rendus de l'Ac.\ d.\ Sc., Déc.\ 1884.

\origpage{242}
above all with several important Memoirs${}^{*)}$, have followed the same direction, either by deepening the general theory, or by studying particular classes of surfaces${}^{**)}$. But the field of research which M.\ \emph{Picard} has opened is far from being exhausted; for the very question of deciding, from the geometric characters of a surface, whether the surface possesses such simple integrals, is not yet resolved in a general manner.

The Italian geometers, on the contrary, have devoted themselves to the study of the linear systems of algebraic curves which belong to an algebraic surface, or, in other terms, to the study of the rational functions of three quantities $x\,y\,z$ bound by an algebraic relation $f(x, y, z) = 0$. This was an order of research which presented itself quite naturally to their attention; for in Italy, during the last years, the two auxiliary theories which serve as the basis of the new developments had been cultivated; namely \emph{the geometry on an algebraic curve}, and \emph{the theory of linear systems of plane curves}.

The geometry on a curve, which took its rise in the immortal works of \emph{Riemann} (1857), was developed later in geometric form in the classic Memoir of MM.\ \emph{Brill} and \emph{Nöther}${}^{***)}$. There yet remained, however, to deliver this theory from the projective considerations which were brought into it too often; and this is what has been done in Italy (\emph{Segre}, \emph{Castelnuovo}...) by reasoning upon the curves without attending either to their orders or to the spaces to which they belong, and by limiting the use of the adjoint curves, whose consideration may sometimes be useful, but is by no means necessary for the development of this theory.

The theory of linear systems of plane curves, from the point of view of the birational transformations of the plane, began with the well known researches of M.\ \emph{Cremona} (1863---1865) and of M.\ \emph{Nöther} (1870), and was developed by the Italian geometers \emph{Caporali}, \emph{Bertini}, \emph{Jung}, \emph{Guccia}, \emph{Martinetti}, \emph{Castelnuovo}... By this route it became possible to approach the study of the properties of the plane which remain invariable with respect to the birational transformations; it is therefore the theory of a particular \emph{class} of surfaces (namely of the \emph{rational} surfaces) that took its rise from these researches.

But it was soon perceived that most of the methods there

\textbf{*)} \emph{Théorie générale des surfaces hyperelliptiques}, Journal de Mathématiques, 4${}^{\text{e}}$ s${}^{\text{e}}$, t.\ IX; \emph{Sur quelques points de la théorie des courbes et des surfaces algébriques}, Journ.\ d.\ Math.\ 4${}^{\text{e}}$ s${}^{\text{e}}$, t.\ X; as well as several Notes in the Comptes Rendus de l'Ac.\ d.\ Sc., 1893---95.

\textbf{**)} Consult also on this subject the Memoir of M.\ \emph{Nöther}, \emph{Ueber die totalen algebraischen Differentialausdrücke}, Math.\ Annalen, 29, 1886.

\textbf{***)} \emph{Ueber die algebraischen Functionen}..., Math.\ Annalen, Bd.\ VII, 1873.

\origpage{243}
employed could have a wider application, for they served also for the study of analogous questions on general algebraic surfaces. One thus came to distinguish two sorts of properties of the linear systems of curves situated on a surface. There are in fact properties which do not depend on the surface, but rather on the nature of the system considered upon it; while there are other properties which belong to all the linear systems of curves situated on the same surface, and depend only on the nature of the latter. It is these last properties that offer the greatest interest, for they furnish characters of a surface which carry over to all the surfaces in birational correspondence with it, \emph{invariant} characters (with respect to the transformations named). It is by the aid of such remarks that the invariants introduced by M.\ \emph{Nöther} were recovered in a new light, and that one succeeded in adding still others to them.

Thus was formed a body of researches closely bound together, although they lie at present scattered through several works published quite recently in Italy${}^{*)}$.

To bring these results together, to present them in the logical order (which is not always the chronological order), to set them beside known properties of curves and of the plane, so as to bring out now the analogies, now the differences: such is the aim we have proposed to ourselves in writing this Monograph. One must not, then, look here for new results, but rather for a new exposition of results already published. This work will nevertheless not appear superfluous to our readers, so we think. The reason is that the original Memoirs which we summarize here do not make easy reading. This depends on several causes. First of all there are the difficulties which come from the theory itself; most of the properties of surfaces are very far from presenting the characters of simplicity which one finds in analogous propositions relative to curves. But the comparison with the latter is very useful for bringing to light the spirit of the new

\textbf{*)} Here, in chronological order, are the principal works whose results we set out here, so far as concerns the general theory of algebraic surfaces; (for the theory of rational surfaces we shall give the relevant citations elsewhere):

\emph{Enriques}, \emph{Ricerche di geometria sulle superficie algebriche}, Mem.\ dell'Acc.\ d.\ Scienze di Torino s${}^{\text{e}}$ 2${}^{\text{a}}$, t.\ 44, 1893.

\emph{Enriques}, \emph{Introduzione alla Geometria sopra le superficie algebriche}, Mem.\ della Società Italiana d.\ Scienze s${}^{\text{e}}$ III, t.\ X, 1896.

\emph{Castelnuovo}, \emph{Alcuni risultati sui sistemi lineari di curve appartenenti ad una superficie algebrica}, ibid., 1896.

\emph{Castelnuovo}, \emph{Sulle superficie di genere zero}, ibid., 1896.

\emph{Enriques}, \emph{Sui piani doppi di genere uno}, ibid., 1896.

\origpage{244}
results. We have recourse to this means unceasingly, and indeed, to make it easier, we have set out in advance the fundamental propositions of the geometry on a curve, in the order which seemed to us preferable.

We thus obtain the advantage of setting before our readers all the notions necessary for approaching the study of surfaces. The reading of this Monograph will therefore require, we think, no more than the knowledge of those general notions of geometry which are familiar even to those who have devoted themselves to quite another order of research.

But to attain this end we have confined ourselves to presenting only the statements of the theorems, and the considerations which serve best to bring them to light; as for the proofs, we have thought it better to suppress them altogether, referring the reader who desires to know them to the original works. The reason is that one does not succeed, for the moment, in setting out these proofs in a manner that avoids every difficulty. The conception which inspires these arguments may well be simple in most cases; but the development requires that one envisage so many particular cases, that one put aside so many objections, that the reader, wearied by this fragmentation of the proof, loses sight of the whole of the results. It is a defect with which theories in formation may ordinarily be reproached; and very often progress consists in the introduction of certain ampler conceptions, of certain more suitable locutions, which allow a great number of cases that seemed distinct to be embraced at a single glance.

The moment has not yet come for introducing such improvements into the theory of surfaces. What matters at present is to surmount the obstacles which one meets on the way, is to reach the summit which shines before us. If there be anyone who will not venture to follow us along a disagreeable route, well then, let them wait! Meanwhile, we hope, they will read with interest the description of the landscape which offers itself to our gaze from the point we have reached; they will find it in this Monograph, which is devoted to it.

But if there are courageous ones who desire to follow us, they will find in the following pages all the information for doing so. And they will be rewarded for their pains, for the field of research is very far from being exhausted, and the region which remains to be explored is still very vast, and may lead to brilliant discoveries. It must be remembered, however, that it is not from a single point of view that so vast a horizon can be embraced. It is on the contrary by the united efforts of several workers following different routes, it is by the learned application of all the resources which geometry and

\origpage{245}
analysis place today at our disposal, that one may hope to deepen the theory of surfaces, and to enrich it with new discoveries.

Only then will this branch be worthy to figure beside the admirable theory of the algebraic functions of one variable, whose construction is a glory of the geometers of our century.

\section*{Chapter I. Birational transformations. Genus of a curve and of a surface after Clebsch and M.\ Nöther.}

\subsection*{1. Birational transformations.}

Projective geometry (the theory of linear and homogeneous substitutions between two groups of variables) has found its most natural extension in the theory of birational transformations between two algebraic varieties.

It is well known how these transformations are defined. Let $x_{1}, x_{2}, \ldots , x_{r}$, and $y_{1}, y_{2}, \ldots , y_{s}$ be two groups of variables, which we may regard as the homogeneous coordinates of a point $x$ in a linear space $X$ of $r-1$ dimensions, and of a point $y$ in a space $Y$ of $s-1$ dimensions; let us establish between these variables relations of the form
\[ \varrho\, y_{i} = f_{i}(x_{1}, x_{2}, \ldots , x_{r}) \qquad (i = 1, 2, \ldots , s) \tag{1} \]
where the $f_{i}$ are algebraic forms all of the same degree in the $x$, and $\varrho$ is a factor of proportionality; one will then have a \emph{rational transformation} between the space $X$ and a certain algebraic variety of the space $Y$. To every general point $x$ of $X$ there corresponds one (single) point $y$ of the space $Y$; while $x$ describes in its space an algebraic line, or surface... $F$, the point $y$ describes in $Y$ (in general) an algebraic line, or surface... $F'$ which corresponds to $F$. Every general point $y$ of $F'$ may, however, come from one, or from several (even from infinitely many) points $x$. If it is the first case that occurs, then, by the aid of the (1) and of the equations which define the $F$ in $X$, one may in turn express the variables $x$ as rational functions of the $y$, by relations of the form
\[ \sigma\, x_{k} = \varphi_{k}(y_{1}, y_{2}, \ldots , y_{s}) \qquad (k = 1, 2, \ldots , r). \tag{2} \]

The (1) and (2) define a \emph{correspondence}, or \emph{birational transformation} between $F$ and $F'$.

\origpage{246}

In particular, one may have a birational (or \emph{Cremona}) transformation between two linear spaces (having the same dimensions); the (1) and (2) define such a transformation between $X$ and $Y$ if $r = s$, and if moreover the (1) can be deduced from the (2), and these from those, for \emph{all} the values of the $x$ and of the $y$.

\subsection*{2. Exceptional elements of a transformation.}

Returning to the varieties $F$ and $F'$ in birational correspondence, there are two cases to be distinguished, according as the $F$, $F'$ are curves, or else varieties of more than one dimension. In the first case, in fact, one has between the curves $F$, $F'$ a one-to-one correspondence \emph{without exceptions} of an essential nature. But if $F$, $F'$ are for example surfaces (of two dimensions), there may well exist on $F$ (or $F'$) simple points (finite in number), to each of which there will correspond on $F'$ (or $F$) all the points of a curve${}^{*)}$. On a surface, a curve which by a suitable birational transformation can be changed into a simple point of another surface is distinguished, by many properties, from the other curves of the first surface; we shall call it an \emph{exceptional curve}. The exceptional curves (in most cases) do not differ from the \emph{ausgezeichnete Curven} considered by M.\ \emph{Nöther}${}^{**)}$; this is what we shall see later.

Suppose in particular that to the point $P$ of $F$ there corresponds the exceptional curve $P'$ of $F'$; and let $C$ be any curve of $F$ passing through $P$. The points of $C$ other than $P$ correspond to the points of a curve $C'$ of $F'$; one might truly regard $C'$ as the curve corresponding to $C$ on $F$. But ordinarily it is convenient to join to the curve $C'$ the exceptional curve $P'$ (which comes from the point $P$ of $C$), and to consider the curve thus composed as the transform of $C$. This is a convention which we shall respect, in general, in what follows.

\subsection*{3. Distribution of algebraic varieties into classes.}

With respect to the birational transformations, algebraic varieties having one and the same number of dimensions are distributed into classes; we shall say, with \emph{Riemann}, that two varieties belong

\textbf{*)} The stereographic projection of a surface of the second order upon a plane gives a very simple example of such curves.

\textbf{**)} Mem.\ cited, Mathem.\ Annalen 8, pag.\ 521.

\origpage{247}
to the same class when a birational correspondence can be established between them. One will thus have to consider, for example, the class of the \emph{rational curves} (or unicursal curves) which can be transformed birationally into a straight line, the class of the \emph{rational surfaces} (or unicursal, homaloidal surfaces) which are in birational correspondence with a plane, etc.

The varieties of one and the same class differ among themselves, it is true, by their projective characters (order, dimension of the space to which they belong,...); but they have several common properties, whose investigation forms the subject of the \emph{Geometry on algebraic varieties}; this branch is therefore occupied with the properties invariable with respect to the birational transformations. It follows that when one studies a \emph{class} of varieties, one may always refer to an abstract, general variety of the class, without attending at all to its projective characters; one might even say that one considers the \emph{algebraic field} of which the varieties of the class (curves, surfaces...) are \emph{projective images}; there would, however, be nothing there but a new word, not a new conception${}^{*)}$.

The consideration of the general variety of a class has this great advantage, that it permits one to abstract from all the singularities of projective nature which a particular variety may present. Thus, in what follows, when we speak of the points of the general variety of a class, we shall always suppose that \emph{simple} points are in question; it is not of course excluded that these points may correspond to multiple points on particular varieties of the class; but one simply affirms that a multiple point is a projective peculiarity of some representative of the class, not of the class itself; it is, so to say, a defect of the image, not of the model.

We shall nevertheless find, in what follows, some paragraphs in which a variety really given by its projective characters is considered (a \emph{plane} curve with its multiple points, a surface \emph{of ordinary space} of known order, etc.); this is because these paragraphs do not properly belong to our theory, but to auxiliary theories. One might even suppress them, if one did not dread the complication which would follow.

\subsection*{4. Genus of an algebraic curve and surface.}

In order to effect the distribution into classes of which we have just spoken, it is necessary to know numerical characters which keep the same value

\textbf{*)} From this point of view one may regard the study of a class of varieties as the study of a \emph{field of rationality} (after \emph{Kronecker}) or of a \emph{field of algebraic functions} (after \emph{Dedekind} and \emph{Weber}).

\origpage{248}
for all the varieties of one and the same class. It is well known how \emph{Riemann}, in studying from the point of view of the \emph{Analysis situs} the real surface which represents the points (real and imaginary) of an \emph{algebraic curve}, arrived at the most important of these characters, which he named the \emph{genus of the curve} (or of the corresponding algebraic relation)${}^{*)}$. \emph{Clebsch} and \emph{Gordan}, in setting out (1866) the same conception in geometric form, gave of the genus the well known definition which follows.

Let $C$ be a plane curve of order $n$; suppose for simplicity that $C$ has no other multiple points besides $d$ double points. Consider the curves of order $n-3$ which pass simply through the $d$ points named, curves which are called \emph{adjoint} to $C$; \emph{the number of adjoint curves which are linearly distinct is the genus $p$ of $C$.}

Beside this \emph{geometric} definition one may give a \emph{numerical} definition of $p$. It suffices to remark that the number of parameters which enter in a homogeneous manner into the equation of a curve of order $n-3$ constrained to pass through $d$ given points is
\[
p = \frac{(n-1)(n-2)}{2} - d.
\]

It is really supposed here that the $d$ double points of $C$ present as many \emph{distinct} conditions to the curves of order $n-3$ which one constrains to pass through them; this is not evident \emph{a priori}, but it was proved later (\emph{Brill} and \emph{Nöther}${}^{**)}$) that the supposition is always verified when one confines oneself to irreducible curves.

Even if one abstracts from this remark, it is well to observe that the two definitions lead equally to invariants of a plane curve; in fact several proofs of the invariance of the genus have been given, some of which (\emph{Riemann}, \emph{Clebsch} and \emph{Gordan},...) relate to the first definition, others (\emph{Cremona}, \emph{Bertini}, \emph{Zeuthen},...) to the second.

The extension of these definitions to algebraic surfaces is immediate in the simplest cases; \emph{Clebsch}${}^{***)}$ and M.\ \emph{Nöther}${}^{\dagger)}$ have indicated it.

Let $F$ be an algebraic surface of order $n$ of ordinary space; suppose that $F$ has no other singularities than a double curve of order $d(\geqq 0)$ and genus $\pi$, and a certain finite number $t(\geqq 0)$ of points triple for the surface, and triple also for the double curve named. Consider now the surfaces of order $n-4$ which pass simply through the double curve (and consequently twice, in general, through its triple points), surfaces which we shall call \emph{adjoint} to $F$; \emph{the}

\textbf{*)} One might trace the notion of the genus back to the theorem of \emph{Abel}, although \emph{Abel} himself does not speak expressly of invariant characters.

\textbf{**)} \emph{Ueber die algebraischen Functionen}..., Math.\ Annalen Bd.\ 7, 1873.

\textbf{***)} Comptes Rendus de l'Acad.\ d.\ Sc., Déc.\ 1868.

\textbf{†)} Memoir cited, \emph{Zur Theorie}...; Mathem.\ Annalen 2, 8.

\origpage{249}
\emph{number of these linearly distinct surfaces} is (according to the proof of M.\ \emph{Nöther}) an invariant character (with respect to the birational transformations) of the surface $F$; it is called the \emph{geometric genus} of $F$, and is designated by $p_{g}$.

Now if we remark that a surface of order $m$ \emph{sufficiently high} must satisfy
\[ md - 2t - \pi + 1 \]
conditions in order to contain the double curve of $F$, we are led (with \emph{Cayley}${}^{*)}$) to form the expression
\[ p_{n} = \frac{(n-1)(n-2)(n-3)}{6} - (n-4)d + 2t + \pi - 1, \]
whose value is exactly $p_{g}$, if $m = n - 4$ is high enough for the preceding formula to be applied. Whether the last condition be fulfilled or not, \emph{the expression of $p_{n}$ always has the property of invariance;} this is what MM.\ \emph{Zeuthen} and \emph{Nöther}${}^{**)}$ have shown. The same expression is perfectly analogous to that considered above in speaking of plane curves; it has been given the name of \emph{numerical genus} of the surface $F$.

But here the analogy with plane curves stops; in fact we can no longer affirm the equality between the values of $p_{g}$ and $p_{n}$; on the contrary \emph{it may well happen that one has $p_{g} \neq p_{n}$ and precisely $p_{g} > p_{n}$.} For example a ruled surface having its plane sections of genus $p$ has (according to \emph{Cayley}${}^{***)}$) $p_{g} = 0$, $p_{n} = -p$. Other examples of surfaces with $p_{g} > p_{n}$ have been brought forward since; we shall return to them in what follows.

\subsection*{5. Remarks on the preceding definitions.}

M.\ \emph{Nöther}, in the Memoir cited several times, advances in the study of surfaces by considering the curves cut out on $F$ by the adjoint surfaces of order $n - 4$. We shall not follow him, for the moment, in this direction; we prefer to show some difficulties to which the definitions of the genera of a surface lead.

In recalling them we have confined ourselves to considering surfaces having quite simple singularities; one might indeed envisage more complicated cases as well, but one is far from being always able to specify in a direct manner the behaviour

\textbf{*)} \emph{On the Deficiency of certain Surfaces}, Mathem.\ Annalen, 3, 1871.

\textbf{**)} \emph{Zeuthen}, \emph{Etude géométrique}..., Mathemat.\ Annalen, 4, 1871; \emph{Nöther}, \emph{Zur Theorie}..., Math.\ Annalen, 8, pag.\ 526.

\textbf{***)} Mem.\ cited.

\origpage{250}
of an adjoint surface at a singular point of the given surface${}^{*)}$. One is therefore not in a position to define the numerical genus of a surface having higher singularities, nor consequently to show in that case the property of invariance of $p_{n}$. It would be far better if one could give of $p_{n}$, and even of $p_{g}$, such a definition as should not depend on the singularities of the surface considered, and should bring out at once the invariance of these characters, and the link between them. This is what we propose to obtain. But for greater clarity it is fitting to start from the simplest notions, and to show how the theory of algebraic curves and surfaces may be reconstructed from a single point of view.

\section*{Chapter II. Linear series of groups on a curve. Linear system of curves on a surface.}

\subsection*{6. Linear series on a curve.}

In the study of a \emph{class} of algebraic varieties the greatest interest belongs to the \emph{rational functions} of the equations which define the class.

Confining ourselves for the moment to curves, if
\[ f(x, y) = 0 \]
is the equation of any plane curve of the class, to every \emph{rational function} $R(x, y)$ of $x$ and $y$ (bound by the $f = 0$) there corresponds a \emph{linear series} $\infty^{1}$ of groups of points on the curve; each group of the series is composed of the points (fixed in number) where the $R$ takes an arbitrarily given value, which changes from one group to another. The number $n$ of the points of each group is called \emph{the order} of the $R$ or of the series; the latter is ordinarily designated by $g_{n}^{1}$. One may also say that the groups of $g_{n}^{1}$ are the groups of the zeros of the rational functions represented by the expression $R(x, y) + C$, where $C$ is a parameter; all these functions have the same poles. More generally let us consider several rational functions of our curve,
\[ R_{1}(x, y), R_{2}(x, y), \ldots , R_{r}(x, y) \]

\textbf{*)} The analytic method proposed to this end by M.\ \emph{Nöther} does not furnish a direct definition permitting the calculation of $p_{n}$, for one must have recourse to a transformation of the surface. See on this subject \emph{Nöther}, Göttinger Nachrichten, 1871, Mathem.\ Annalen, 29, as well as \emph{Zeuthen}, \emph{Révision et extension des formules numériques}... Mathem.\ Annalen 10, pag.\ 544, 1876.

\origpage{251}
all having the same $n$ poles, and not being bound by any linear relation with constant coefficients; let us form the function
\[ \lambda_{1} R_{1} + \lambda_{2} R_{2} + \cdots + \lambda_{r} R_{r} + C \]
which depends on the parameters $\lambda_{1}, \ldots , \lambda_{r}, C$; the group of the zeros of this describes (when the parameters are made to vary) \emph{a linear series $g_{n}^{r}$ of order $n$, and dimension $r$;} $r$ general points of the curve enter into one group of $g_{n}^{r}$, and into one only. The group of the poles also belongs to the series, and presents no peculiarity in it. It is now easy to define the widest linear series of order $n$ to which a group of $n$ given points belongs, or, as we shall say, the \emph{complete series} (\emph{Vollschaar}) which is determined by that group. It suffices in fact to consider \emph{all} the rational functions of order $n$ which have their poles in that group; the groups of the zeros of these functions form the series required. It is clear that the dimension of the last series cannot exceed $n$; if it attains this value, the curve in question is rational, for in the contrary case one could not fix arbitrarily the zeros and the poles of one and the same rational function.

The same considerations may also be presented in geometric form${}^{*)}$. Let $C$ be an algebraic curve (plane, or of ordinary space, or of a hyperspace); let us cut it by a linear system $\infty^{r}$ $(r \geqq 0)$ of varieties (curves, surfaces...)
\[ \lambda_{1} f_{1} + \lambda_{2} f_{2} + \cdots + \lambda_{r+1} f_{r+1} = 0, \]
where the $f$ are algebraic forms all of the same degree (in three, four... variables). If each $f$ has $n$ points of intersection with the $C$, and if there is in the system no variety containing the $C$, the groups of the intersections will form a linear series $g_{n}^{r}$. More generally, if with the groups of the $g_{n}^{r}$ we form groups of $n \pm n'$ points, by adding to the first groups the same $n'$ points fixed arbitrarily on $C$, or by subtracting them if they enter into every group of $g_{n}^{r}$, we shall say that the new groups form a linear series $g_{n \pm n'}^{r}$.

The varieties $f$ which have allowed us to construct the series $g_{n}^{r}$ are not, however, bound up with it (from the point of view of the geometry on the curve), and do not enter at all into the study of the series. One may

\textbf{*)} As we have said in the introduction, we present here the geometric theory of the linear series (which is due to MM.\ \emph{Brill} and \emph{Nöther}, Mem.\ cited) following the plan adopted in some Italian works. The developments and the relevant citations will be found in the Monograph of M.\ \emph{Segre}, \emph{Introduzione alla geometria sopra un ente algebrico}..., Annali di Matematica s${}^{\text{e}}$.\ II, tomo XXII, 1894. In the same volume one may also consult with profit a Monograph of M.\ \emph{Bertini}, where the same subject is set out from the point of view of \emph{Brill} and \emph{Nöther}.

\origpage{252}
indeed construct the latter by other systems of varieties, whether on the same curve $C$, or on a curve $C'$ of the same class (for a birational transformation between $C$ and $C'$ changes the $g_{n}^{r}$ of $C$ into a $g_{n}^{r}$ of $C'$).

One may even define the linear series $g_{n}^{r}$ on any algebraic curve without expressly mentioning the varieties which serve to cut it out. In fact \emph{a $g_{n}^{r}$ is every system $\infty^{r}$ of groups of $n$ points on the curve whose groups depend rationally on $r$ parameters, in such a manner that the condition of containing a point arbitrarily fixed on the curve translates into a linear relation between the parameters.} It follows naturally that $r$ general points of the curve define a group; but we shall see presently how one may also (with some restrictions) draw the premises from the conclusion.

When a linear series $g_{n}^{r}$ $(r \geqq 0)$ is given on a curve, it may be that there exists on the same curve a new series $g_{n}^{s}$ of the same order and of dimension $s > r$, containing all the groups of $g_{n}^{r}$, and consequently the $g_{n}^{r}$ itself; it may on the contrary happen that such a series $g_{n}^{s}$ does not exist. Now one says \emph{that a series is complete when it is not contained in another series of the same order and of higher dimension;} the series is \emph{incomplete} in the contrary case${}^{*)}$.

In this last case, by forming series of the order $n$ wider and wider, which contain a given incomplete series $g_{n}^{r}$, one will necessarily arrive at a complete series $g_{n}^{s}$ $(r < s \leqq n)$ containing the former. Now it is proved that this last series does not depend on the procedure followed in constructing it; it depends only on the series $g_{n}^{r}$ from which one starts. In other terms:

\emph{A linear series $g_{n}^{r}$ $(r \geqq 0)$ either is complete, or else is contained in a complete series of the same order $n$, which is entirely determined by the $g_{n}^{r}$}${}^{**)}$.

It follows for example that \emph{if from a complete series $g_{n}^{r}$ one deduces a new series $g_{n-\nu}^{\varrho}$ by imposing on the groups of the former the condition of containing $\nu$ arbitrary points of the curve, and by then suppressing them in those groups, the new series will again be complete.} It is called the \emph{residual series} of the $\nu$ points with respect to $g_{n}^{r}$; its dimension $\varrho$ is equal to $r - \nu$, or greater, according as the $\nu$ points present conditions all independent, or not, to the groups of $g_{n}^{r}$ which contain them.

It is the complete series that offer the greatest interest. The theory of surfaces will nevertheless often lead us to consider incomplete series as well. A character of such a series $g_{n}^{r}$ to which regard must be had is the difference $s - r$ which one obtains by

\textbf{*)} One sees clearly that this definition agrees perfectly with that which we have given, starting from analytic considerations.

\textbf{**)} \emph{Segre}, l.\ c., §~14.

\origpage{253}
subtracting the dimension of our series from the dimension of the complete series $g_{n}^{s}$ which contains the former. We call this difference the \emph{defect (deficienza, difetto di completezza, Defect)} of the $g_{n}^{r}$.

\emph{Remark.} --- On the definition of a linear series a remark may be made which, for the rest, has no importance for the sequel of our Monograph. Let us call an \emph{involution of order $n$ and dimension $r$} on an algebraic curve every algebraic system of $\infty^{r}$ groups of $n$ points such that $r$ general points of the curve enter into a single group of the system. Then we may say that every linear series is an involution; will the converse theorem also be true? Not always, for there are involutions $\infty^{1}$ which are not linear; those, for example, which are cut out by the generators of an irrational ruled surface on the curves belonging to the surface. But, \emph{apart from} the involutions $\infty^{1}$, and \emph{from the involutions $\infty^{r}$ every group of which is formed by $r$ arbitrary groups of an involution $\infty^{1}$} (of order $n \geqq 1$), it is proved that \emph{every involution of dimension $> 1$ is a linear series}${}^{*)}$. One has there, then, a new definition of these series.

\subsection*{7. Linear system of curves on a surface.}

The rational functions of an algebraic equation in two variables have led us to consider the linear series on a curve; likewise, starting from an equation in three variables, one would come to consider the linear systems of curves on a surface.

To define them from the geometric point of view it suffices to consider an algebraic surface $F$ (of ordinary space or of a hyperspace) and a linear system of varieties (of two, ..., dimensions)
\[
\lambda_{1}f_{1} + \cdots + \lambda_{r+1}f_{r+1} = 0;
\]
the curves which these varieties cut out on the $F$ form a \emph{linear system} whose \emph{dimension} is $r$, if none of the varieties named contains the $F$.

As regards the varieties $f$, it may be that they contain one and the same fixed curve of $F$, a curve which one is free to regard as forming part of the general curve of the system, or else as foreign to it (which will then reduce to the variable part of the intersection $fF$). It may further happen that the $f$ pass through fixed points of $F$ (apart from the fixed curves); these points, finite in number, which belong to all the curves of the system, are called its \emph{base points}.

\textbf{*)} This theorem was given at the same time (June 1893) by MM.\ \emph{Humbert} (Comptes Rendus de l'Ac.\ d.\ Sc.; see also Journal de Mathématiques, 4${}^{\text{e}}$ série t.\ X) and \emph{Castelnuovo} (Atti dell'Accad.\ d.\ Sc.\ di Torino).

\origpage{254}

A linear system $\infty^{r}$ of curves on a surface has the property that through $r$ general points of the surface there passes a single curve of the system. This property may even serve as a definition for a linear system of dimension $r > 1$. There exist, it is true, non-linear systems $\infty^{1}$ of curves satisfying the condition that through every general point of the surface there passes a single curve of the system; \emph{these are the irrational pencils}, of which the system of the generators of an irrational ruled surface furnishes a very simple example. \emph{But every system $\infty^{r}$, with $r > 1$, a single curve of which is entirely determined by $r$ points of the surface, is a linear system, if one excludes only the case where the general curve of the system is composed of several curves of an irrational pencil}${}^{*)}$.

The general curve of a linear system may be \emph{irreducible}, or else be composed of several parts; in the first case the system is said to be \emph{irreducible}, and \emph{reducible} in the second case. The study of this last case reduces to the first (the more interesting) by the aid of the theorem:

\emph{The general curve of a reducible linear system $\infty^{r}$ is formed:}

1) \emph{either by a fixed part (common to all the curves of the system), joined to a curve which varies in an irreducible linear system $\infty^{r}$;}

2) \emph{or else by several $(\geqq r)$ curves which belong to one and the same pencil (rational or not), to which fixed parts may sometimes be added.}

The exclusion of the reducible systems, which would render our considerations simpler, is not always possible; for there are operations on linear systems which, even starting from an irreducible system, may lead to reducible systems. But, by the aid of suitable conventions, most of the properties of the one and the other systems may be gathered under a single statement.

Let us confine ourselves for the moment to irreducible systems; let $|C|$ be such a system, $C$ its general curve.

The most important characters of the system (which remain invariable under the birational transformations of the surface) are the following:

1) the \emph{dimension} $r$ of $|C|$;

2) the \emph{degree} $n$, that is to say the number of variable intersections of two general curves of the system (for $r = 1$, one has $n = 0$).

3) the \emph{genus} $\pi$, which is the genus of the curve $C$.

One has also often to consider a linear series which is invariably bound up with the system $|C|$, and which will play in what follows an

\textbf{*)} \emph{Enriques}, \emph{Una questione sulla linearità}..., Rendic.\ della R.\ Accad.\ d.\ Lincei, luglio 1893; the proof is reported with some simplifications in the cited Monograph of M.\ \emph{Segre}, n${}^{\text{o}}$ 27.

\origpage{255}
important role; it is the \emph{characteristic series $g_{n}^{r-1}$}, which is cut out on a $C$ by the other curves of $|C|$${}^{*)}$. Here one has $r - 1 \leqq n$, whence one concludes that the dimension of an irreducible system cannot exceed the degree increased by one unit.

\emph{Remark.} --- The property of a system $|C|$, situated on a surface $F$, of being reducible or irreducible, as well as the characters $r$, $n$, $\pi$ of $|C|$, carry over without modification to a system $|C'|$ of a surface $F'$, when there is a birational transformation which changes $F$ into $F'$ and $|C|$ into $|C'|$. In order, however, that this remark should subsist without exceptions, the following convention must be made. It may happen that a base point $P$ of $|C|$ on $F$, having the order $i \geqq 1$ for the general curve $C$ (a point which we shall suppose simple for $F$), gives, by the transformation, an exceptional curve $P'$ of $F'$. According to n${}^{\text{o}}$ 2, one ought to regard this curve $P'$ as forming part of the curve $C'$ transformed from $C$; but then $|C'|$ would be reducible, even though $|C|$ is not. To avoid the discrepancy it is agreed in this case to abstract from the curve $P'$ (corresponding to a base point of $|C|$), and to regard $C'$ as the locus of the points corresponding to the points of $C$, except the point $P$. On the contrary the rule of n${}^{\text{o}}$ 2 will be respected when $P$ (which is transformed into the curve $P'$) is not a base point of $|C|$, but a general point of $F$.

Besides the example we have brought forward, there would be others analogous to be discussed. We do not insist upon them here; the reader who wishes to go deeper into these questions will be able to do so elsewhere${}^{**)}$; they will see, however, from these few words, that there are in the theory of surfaces difficulties and distinctions which one does not succeed in avoiding.

\subsection*{8. Complete linear systems.}

If we propose to carry over to the linear systems on a surface the notion of complete series established for a curve, we meet first of all a distinction which did not present itself there. In fact, if one seeks a system $|D|$ containing the given system $|C|$, it is necessary to settle which character of $|C|$ is to be transmitted to the new system. One may require that the genus $\pi$ alone should remain

\textbf{*)} The importance of this series in the study of linear systems of plane curves was remarked by MM.\ \emph{Segre} (\emph{Sui sistemi lineari}..., Rendic.\ Circolo Mat.\ di Palermo, tomo I, 1887) and \emph{Castelnuovo} (\emph{Ricerche generali sopra i sistemi lineari}..., Memorie dell'Acc.\ d.\ Sc.\ di Torino, s${}^{\text{e}}$ II, t.\ 42). As for the consideration of this series on surfaces one has only to consult the works of MM.\ \emph{Enriques} and \emph{Castelnuovo}, cited in the introduction.

\textbf{**)} \emph{Enriques}, \emph{Introduzione}..., cap.\ I.

\origpage{256}
invariable; but one may also require that it be the degree $n$ (and consequently the genus) which does not change. Although the first condition is of a more general nature than the second, it is to the latter that more interest must be assigned; the progress of the theory of surfaces has shown it. Consequently we shall say that \emph{an irreducible linear system of curves is complete (with respect to the degree) when there is on the surface no other system of the same degree and of higher dimension which contains the first system}${}^{*)}$.

Here is a very simple example of it. Consider on a plane the linear system of all the curves having given multiplicities at certain base points; it is a \emph{complete} system. In other terms: \emph{a linear system of plane curves entirely defined by its base points is complete}${}^{**)}$.

The definition which we have just given above may be presented in a more general manner, which permits it to be extended to the reducible systems as well. To this end the locution must first be explained: \emph{a curve $C_{0}$ (simple or composite) is contained totally in a given system $|C|$.}

This amounts to saying that the curve $C_{0}$ is by itself \emph{a} curve of $|C|$ (it is thus excluded that it should be \emph{a part} of such a curve), and moreover that $C_{0}$ has the same order of multiplicity as the general curve of $|C|$ at every base point of the system. Similarly a system is said to be contained \emph{totally} in another when the general curve of the first system is contained totally in the second, and the base points of the latter are also the base points of the former. It is then easily proved that two systems of which one is contained totally in the other have the same \emph{degree}; and \emph{viceversa}, if of two systems having the same degree one is contained in the other, one may affirm that it is contained in it \emph{totally}${}^{***)}$. It results that the definition of a complete system may now be stated as follows:

\emph{A linear system is complete if it is not contained totally in another system.}

Under this form the definition has a meaning, and shows what

\textbf{*)} The notion of complete system (with respect to the degree), \emph{sistema normale}, was introduced by M.\ \emph{Enriques} (\emph{Ricerche di geometria}... I, 2; \emph{Introduzione}...): the extension to the varieties $\infty^{r}$ is found in the cited Monograph of M.\ \emph{Segre} (n${}^{\text{o}}$ 26). M.\ \emph{Enriques} in the first Memoir speaks also of the systems complete with respect to the genus, which he calls \emph{sistemi completi}. We do not have recourse to them in this Memoir.

\textbf{**)} If among the base points there are simple ones, by abstracting from these one obtains a system wider than the given system, which has nevertheless the same genus as the latter, but a higher degree. One sees therefore that a system complete with respect to the degree may well not be complete with respect to the genus.

\textbf{***)} \emph{Enriques}, \emph{Introduzione}... n${}^{\text{o}}$ 7.

\origpage{257}
is to be understood by \emph{complete system}, even when the given system is reducible, or if it has the dimension \emph{one} or \emph{zero}; in these last cases, however, one must designate which are the points to be regarded as base points of the complete system.

It is in this sense that the following theorem is to be understood, which is perfectly analogous to a theorem on linear series${}^{*)}$:

\emph{Every linear system $\infty^{r}$ $(r \geqq 0)$ is complete, or else it is contained totally in a complete system which is entirely determined by it.}

One has again this other theorem, which also recalls a theorem on series:

\emph{If from a complete linear system one deduces a new system by imposing on the curves of the first the condition of passing through fixed points with given orders of multiplicity, the last system is also complete.}

In what follows we shall consider, as far as possible, complete systems.

\subsection*{9. Elementary operations on series and on systems of curves.}

We call by this name the operations of addition and of subtraction of series and of systems, which we are about to define; they correspond to the multiplication and the division of the rational functions which belong to one and the same field of rationality.

Let there be first two complete linear series $g_{m}^{q}$, $g_{n}^{r}$ on one and the same curve; let us form all the groups of $m + n = p$ points which one obtains by uniting an arbitrary group of $g_{m}^{q}$ to an arbitrary group of $g_{n}^{r}$; the groups thus formed belong, as is easily proved, to one and the same complete series $g_{p}^{s}$, which is called the \emph{sum of the two given series}. The case where the two series $g_{m}^{q}$, $g_{n}^{r}$ coincide is not excluded; one arrives then at the \emph{double} series of a given series. And in the same way one defines the series sum of several series, and the multiple (k.\ uple) series of a given series${}^{**)}$.

We are now in a position to define the inverse operation, the subtraction of series, so far as it is possible. Let us consider to

\textbf{*)} \emph{Enriques}, \emph{Introduzione}... n${}^{\text{o}}$ 11.

\textbf{**)} \emph{Segre}, \emph{Introduzione}... n${}^{\text{o}}$ 59. Truly there was sometimes understood by \emph{sum of two series} $g_{m}^{q}$, $g_{n}^{r}$ (complete, or not) the series of order $m + n$ and of the smallest dimension which contains the groups composed of a group of $g_{m}^{q}$ and of a group of $g_{n}^{r}$ (\emph{see} \emph{Castelnuovo}, \emph{Sui multipli di una serie lineare}..., Rendic.\ Circolo Mat.\ di Palermo, t.\ VII). But in this Monograph we adopt the definition set out above.

\origpage{258}
this end two complete series $g_{p}^{s}$, $g_{n}^{r}$ $(p > n)$ on the curve, and suppose that a certain group $G_{n}$ of the second series is contained in $\infty^{q}$ groups $(q \geqq 0)$ of the first. The groups of the $p - n = m$ points which remain form, as we know, a complete series $g_{m}^{q}$ \emph{residual} to $G_{n}$ with respect to $g_{p}^{s}$. Now this series $g_{m}^{q}$ has a remarkable property; in fact it is residual (with respect to $g_{p}^{s}$) not only to $G_{n}$, but also to every other group of $g_{n}^{r}$. One says consequently that $g_{m}^{q}$ \emph{is the residual series of} $g_{n}^{r}$ \emph{with respect to} $g_{p}^{s}$; but one sees likewise that $g_{n}^{r}$, in its turn, is residual to $g_{m}^{q}$ with respect to $g_{p}^{s}$; the last series is the sum of the two first. This property gives rise to the statement which follows${}^{*)}$:

\emph{Two complete series $g_{p}^{s}$, $g_{n}^{r}$ being given on one and the same curve, if a group of the $g_{n}^{r}$ is contained in some group of the first series, it is the same for every other group of the $g_{n}^{r}$; there exists then a third series $g_{m}^{q}$ entirely determined, which added to the second $g_{n}^{r}$ gives the first series $g_{p}^{s}$; the two series $g_{m}^{q}$, $g_{n}^{r}$ are residual to one another with respect to $g_{p}^{s}$.}

To this theorem one might give the name of \emph{theorem of the remainder} for series; this name has, however, already been adopted by MM.\ \emph{Brill} and \emph{Nöther} to designate a theorem which has much analogy with the preceding, but which comprises also other notions of which we shall speak later.

Let us come now to the linear systems of curves on a surface; the analogy permits us to extend at once the preceding results. If $|C_{1}|$ and $|C_{2}|$ are two complete systems of curves on one and the same surface, the curves $C_{1} + C_{2}$, which may be formed by uniting two curves taken arbitrarily in the two systems, belong to one and the same determinate complete system $|C| = |C_{1} + C_{2}|$, which is called the \emph{sum of} $|C_{1}|$ \emph{and} $|C_{2}|$. In particular one may speak of the \emph{double, triple}... system of a given system.

Let us start now from the system $|C|$, and let $C_{1}$ be a curve which enters into some (reducible) curve of $|C|$; the curves $C_{2}$ which with the $C_{1}$ give rise to curves of $|C|$ form a complete system $|C_{2}|$ \emph{residual to the} $C_{1}$ \emph{with respect to} $|C|$. Now if the $C_{1}$ belongs to a complete system $|C_{1}|$ (determined by it), it is proved that $|C_{2}|$ is also residual (with respect to $|C|$) to every other curve of $|C_{1}|$; it is \emph{the system} $|C_{2}| = |C - C_{1}|$ \emph{residual to} $|C_{1}|$ \emph{with respect to} $|C|$, just as $|C_{1}|$ is residual to $|C_{2}|$ with respect to $|C|$. One has therefore the theorem${}^{**)}$:

\emph{If a complete system $|C|$ contains a curve of a second complete}

\textbf{*)} \emph{Segre}, \emph{Introduzione}... n${}^{\text{o}}$ 58.

\textbf{**)} \emph{Enriques}, \emph{Introduzione}... n${}^{\text{o}}$ 13.

\origpage{259}
\emph{system $|C_{1}|$, the first system contains every curve of the second; there exists then a complete system $|C_{2}|$ entirely determined, such that $|C| = |C_{1} + C_{2}|$; each of the two systems $|C_{1}|$ and $|C_{2}|$ is residual to the other with respect to $|C|$.}

It must be remarked, however, that even if $|C|$ and $|C_{1}|$ are supposed irreducible, it may well happen that $|C_{2}| = |C - C_{1}|$ is reducible.

To conclude these considerations on the system sum of two given systems, we shall make known here two very simple relations, to which one must nevertheless often have recourse in researches on surfaces. These formulas express the degree $n$ and the genus $\pi$ of the sum system $|C| = |C_{1} + C_{2}|$ by the aid of the analogous characters $n_{1}$, $\pi_{1}$ and $n_{2}$, $\pi_{2}$ of the two systems $|C_{1}|$ and $|C_{2}|$; one has precisely
\[
\begin{aligned}
n &= n_{1} + n_{2} + 2i,\\
\pi &= \pi_{1} + \pi_{2} + i - 1,
\end{aligned}
\]
where $i$ designates the number of variable intersections of a $C_{1}$ with a $C_{2}$. It is supposed here that $|C_{1}|$ and $|C_{2}|$ are irreducible; but it is useful to remark that if it were not so, one might likewise define the genus and the degree of the reducible systems in such a way that the preceding relations should still be true${}^{*)}$.

\section*{Chapter III. Curves adjoint to a plane curve. --- Surfaces sub-adjoint to a surface of ordinary space.}

In the preceding Chapter we have considered complete series and complete systems of curves, but we have not indicated the means of \emph{completing} a given series or system. This is what we shall do now, confining ourselves (in this Chapter) to the series situated on \emph{plane} curves, and to the systems of curves on the surfaces of \emph{ordinary space}. All the other cases may be brought back to these by the aid of a suitable birational transformation of the given curve or surface.

\subsection*{10. Curves adjoint to a plane curve.}

Let $C$ be first a \emph{plane} curve of the order $m$, which will in general have multiple points. One gives the name of \emph{adjoint curve} to $C$ to a curve which passes $\alpha - 1$ times through every multiple point of order $\alpha$ of the curve $C$ $(\alpha = 2, 3 \ldots)$${}^{**)}$. The curves adjoint to $C$ of a given order (which

\textbf{*)} \emph{Enriques}, \emph{Introduzione}... n.\ 15, 16.

\textbf{**)} \emph{Brill} and \emph{Nöther}, \emph{Ueber die algebraischen Functionen}..., l.\ c.

\origpage{260}
certainly exist if that order is high enough) form a linear (complete) system of plane curves. Now this system has a very remarkable property, whence follows the application of the adjoint curves to the theory of linear series. In fact${}^{*)}$:

\emph{The curves adjoint to $C$ of any given order cut out on the curve $C$ a complete linear series;} it must be understood here that every group of the series is formed by the intersections of $C$ with an adjoint curve outside the multiple points. It follows (n${}^{\text{o}}$ 9) that one likewise obtains on $C$ a complete series if one obliges the adjoint curves to pass through points fixed arbitrarily \emph{on} $C$, and confines oneself to considering the intersections which fall outside these points (and the multiple points).

By the aid of the last observation one may at once \emph{construct on the curve $C$ the complete linear series which is defined by a group of points, or by a given incomplete series}. Through the given group $G$ (which in the second case will be any group of the series) one makes pass any adjoint curve, which is always possible if the order of the latter is high enough. The curve cuts out on $C$, outside the multiple points and the group $G$, a group of points $G'$. Let us draw through $G'$ all the adjoint curves of the chosen order; these will cut out on $C$ the complete series defined by $G$.

There is something arbitrary in this construction, for one disposes of the order of the adjoint curves, and also of the curve, among these, which one first draws through $G$; but in whatever manner one arranges matters, one will always arrive in the end at the same series. It is precisely this affirmation that constitutes the \emph{Restsatz} of MM.\ \emph{Brill} and \emph{Nöther}${}^{**)}$.

\subsection*{11. Characteristic series of a linear system of plane curves.}

The property of cutting on a given plane curve $C$ a complete series belongs not only to the systems of adjoint curves,

\textbf{*)} \emph{Brill} and \emph{Nöther}, l.\ c.\ pag.\ 275; see also \emph{Segre}, \emph{Introduzione}... n${}^{\text{o}}$ 77.

\textbf{**)} The \emph{Restsatz} thus proceeds from two quite distinct propositions, one of which belongs properly to the \emph{Geometry on the curve}, and relates to the operations of addition and subtraction of series (n${}^{\text{o}}$ 9); while the other proposition belongs to the \emph{Geometry of the plane}, and expresses the property of the systems of adjoint curves which we have stated above. One might even remark that the \emph{Restsatz} (as it is stated by MM.\ \emph{Brill} and \emph{Nöther}) says a little less than the union of our two propositions, for it confines itself to affirming that the series cut out by the adjoint curves in the manner indicated does not depend on these, while we affirm besides that the same series is complete; but it must moreover be remarked that this last part is deduced at once from the first. See \emph{Brill} and \emph{Nöther}, l.\ c.\ pag.\ 273; \emph{Segre}, \emph{Introduzione}..., n${}^{\text{o}}$ 78.

\origpage{261}
but also to other systems. For example the curves of an arbitrary order which pass with the multiplicity $\alpha$ through each multiple point of order $\alpha$ of $C$ themselves cut out on $C$ a complete series; this is a property which follows easily from the property of the adjoint curves. The most interesting case presents itself if the order of the curves we are considering is the order of $C$; in fact in this case the series named is nothing else than the \emph{characteristic series} of the complete system of the plane curves having the order of $C$ and the same multiplicities at the multiple points of the latter. We may therefore state the following theorem, which is going to be useful to us${}^{*)}$:

\emph{The characteristic series of a complete linear system of plane curves is always complete.}

\subsection*{12. Surfaces sub-adjoint to a surface of ordinary space.}

We shall now see how the preceding theorems may be extended to the surfaces of ordinary space.

Let $F$ be such a surface having arbitrary singularities (multiple curves and multiple points). One may always construct a surface $\Psi$ of a sufficiently high order which fulfils the condition of cutting on every general plane a curve adjoint to the section of $F$ made by the same plane. One specifies in this manner the behaviour of the $\Psi$ along the multiple curves of $F$; thus when ordinary multiple curves are in question, the $\Psi$ will pass $\alpha - 1$ times through every curve of $F$ which has the multiplicity $\alpha$ for $F$. Nothing is said, on the contrary, about the passage of $\Psi$ through the \emph{isolated} multiple points of $F$, whose presence on $F$ does not follow (so to say) from the existence of the multiple curves; so that if $F$ possesses, for example, a multiple point belonging to no multiple curve, the $\Psi$ will not, in general, pass through that point. In virtue of these explanations we shall say that the $\Psi$ is a surface \emph{adjoint to $F$ along the multiple curves}, or, for short, that the $\Psi$ is \emph{sub-adjoint (subaggiunta) to $F$}${}^{**)}$. We reserve the name of \emph{adjoint} surfaces for certain surfaces, with which we shall have to occupy ourselves later, which, while satisfying the conditions of the $\Psi$, will further pass suitably through the isolated points of $F$. For the moment the sub-adjoint surfaces suffice for our purpose.

\textbf{*)} \emph{Brill} and \emph{Nöther}, l.\ c.\ pag.\ 308; \emph{Castelnuovo}, \emph{Ricerche generali sopra i sistemi lineari}... l.\ c.\ n${}^{\text{o}}$ 18.

\textbf{**)} \emph{Enriques}, \emph{Introduzione}... n${}^{\text{o}}$ 17; M.\ \emph{Nöther} (\emph{Zur Theorie}..., Math.\ Annalen 8, pag.\ 509) calls \emph{adjungirte Flächen} the surfaces which we call \emph{sub-adjoint}, and reserves the name of \emph{surfaces} $\varphi$ for the surfaces of order $n - 4$ which, according to this Monograph, are \emph{adjoint to $F$}. The development of the theory of surfaces has counselled this change of locutions.

\origpage{262}

The surfaces of a given order sub-adjoint to $F$ form a linear system which enjoys a property analogous to that of the adjoint curves. In fact${}^{*)}$:

\emph{The surfaces of any given order which are sub-adjoint to a surface cut out on the latter a complete system of curves:} the general curve of the system is formed by the intersection of the sub-adjoint surface with $F$, outside the multiple curves of $F$. The theorem still subsists if one obliges the sub-adjoint surfaces to pass through points or curves assigned on $F$, and subtracts these fixed curves from the curves of the system (n${}^{\text{o}}$ 9).

By the aid of these remarks one may always construct on $F$ the complete system determined by a given curve $C$, or by an incomplete linear system. It suffices in fact to draw through the curve $C$ (which in the second case will be chosen arbitrarily in the system) a sub-adjoint surface $\Psi$ of an order high enough for the construction to be possible. The $\Psi$ will in general cut out on $F$ a new curve $C'$, outside $C$ and the multiple curves of $F$; let there now be drawn through $C'$ all the adjoint surfaces of the same order as $\Psi$; these will cut out on $F$ a complete linear system to which the curve $C$ belongs; it is the system sought, or at most it differs from it only in this, that the system sought may have some base points more. It suffices, in this last case, to impose these new base points on the curves of the system we have just constructed, or (what comes to the same) to impose on the surfaces $\Psi$ suitable passages or contacts with $F$ at the same points.

The complete system at which one arrives depends neither on the order of the sub-adjoint surfaces $\Psi$ that are employed, nor on the surface, among the $\Psi$, which is drawn first. This affirmation constitutes the \emph{Restsatz} for surfaces, after M.\ \emph{Nöther}${}^{**)}$.

\emph{Remark.} --- The theorem on the characteristic series of a linear system of plane curves may also be extended to the linear systems of surfaces of ordinary space. Let us state the result, although we shall not in what follows have occasion to have recourse to it:

\emph{If a linear system of surfaces is entirely determined by its base curves and base points, the system of curves which is cut out on one of the surfaces by the other surfaces of the system is complete.}

\textbf{*)} \emph{Enriques}, \emph{Introduzione}... n${}^{\text{o}}$ 35 (there the adjoint surfaces are really considered, but the thing is true also for the sub-adjoint surfaces). The same theorem follows also from the \emph{Restsatz} of M.\ \emph{Nöther} of which we shall speak presently.

\textbf{**)} \emph{Zur Theorie}..., Math.\ Annalen 8, pag.\ 509; there would be occasion to repeat here a remark analogous to that which we have made apropos of plane curves.

\origpage{263}

\section*{Chapter IV. System adjoint to a given system.}

We shall define in this Chapter an operation which is to play a fundamental role in the theory of the linear systems situated on a surface; it is the operation of \emph{adjunction}. It furnishes the means of deducing from every given linear system a new system, \emph{the adjoint} to the first; and by the consideration of the bonds which pass between these two systems it leads us to establish the principal invariants of a surface with respect to the birational transformations.

Following our method, we begin with considerations relative to the linear series on curves.

\subsection*{13. Canonical series on a curve.}

Among the curves adjoint to a plane curve $C$ of order $m$, those which have the order $m - 3$ must be considered in particular. Suppose that such curves exist, and let us designate by $p(>0)$ the number of those which are linearly independent. Then the said curves cut out on the $C$ a complete linear series whose dimension is $p - 1$, and whose order is (according to MM.\ \emph{Brill} and \emph{Nöther}${}^{*)}$) $2p - 2$; it is a $g_{2p-2}^{p-1}$. Now this series enjoys a fundamental property. If in fact one transforms birationally the curve $C$ into another plane curve $C'$ of order $m'$, the series $g_{2p-2}^{p-1}$ of $C$ is transformed into the series which is in its turn cut out on $C'$ by the curves of order $m' - 3$ adjoint to $C'$. In other terms \emph{the series cut out on a plane curve of order $m$ by the adjoint curves of order $m - 3$ enjoys the property of invariance} (with respect to the birational transformations). We shall often have occasion to name this series; we shall designate it in what follows by the name of \emph{canonical series} on the curve $C$ (or on any other curve of the same \emph{class}). One will therefore speak also of the canonical series on a skew curve; it will be the series corresponding to that cut out by the adjoints of order $m - 3$ on a plane curve of order $m$ in birational correspondence with the given curve.

The dimension of the canonical series increased by one unit is also an invariant of the curve considered; it is the \emph{genus} $p$ of the curve, whose definition we have already recalled in a particular case (n${}^{\text{o}}$ 4).

\textbf{*)} \emph{Ueber die algebraischen Functionen}... pag.\ 280; see also \emph{Segre}, \emph{Introduzione}..., n${}^{\text{o}}$ 73.

\origpage{264}

The character $p$, in virtue of its definition, cannot descend below zero; it is null if the corresponding curve, supposed plane, of order $m$, has no adjoint curves of order $m - 3$. It is well known moreover that the condition $p = 0$ entirely defines a \emph{class} of curves; these are the \emph{rational curves}.

This case excepted, there are known properties of the canonical series which it is fitting to recall here, for they are in continual use${}^{*)}$. The canonical series is the only series which has, on a curve of genus $p$, the dimension $p - 1$ with the order $2p - 2$. The same series has no fixed points common to all its groups. Consequently, if one fixes any point of the curve, one deduces from the canonical series $g_{2p-2}^{p-1}$ a new series $g_{2p-3}^{p-2}$; this in its turn has \emph{in general} no fixed points. At most, in a particular case, it may have one fixed point; it then happens that all the groups of the canonical series which contain any fixed point contain in consequence a second fixed point which is determined by the first; there exists then on the curve a series $g_{2}^{1}$ such that any $p - 1$ of its groups form a group of the canonical series $g_{2p-2}^{p-1}$; this is the case of the \emph{hyperelliptic} curves.

\emph{Remark.} --- Of the canonical series one may give a well known analytic definition, which has the advantage of extending at once to all the curves of one and the same class, whether they be plane or skew: \emph{the groups of the canonical series are the groups of the zeros of the Abelian differentials of the first kind which belong to the given curve.}

\subsection*{14. System adjoint to a system of plane curves.}

Let us return to the plane curve $C$ of order $m$, and suppose now that it is a general curve of a linear system $|C|$ of plane curves, a complete system, that is to say one entirely determined by its base points. These points have the same orders of multiplicity $(\geqq 1)$ for the $C$ and for the other curves of the system; and all the multiple points of the $C$ fall among them. It follows that the curves adjoint to $C$, of the order $m - 3$, (and more generally of any order) are likewise adjoint to every other curve of the system $|C|$. We shall say that these curves of the order $m - 3$ form \emph{the system} $|C'|$ \emph{adjoint to the system} $|C|$; one says \emph{adjoint}, for short, without designating the order of the curves $C'$, for one has no occasion to have recourse to adjoint systems of a different order.

\textbf{**)} \emph{Brill} and \emph{Nöther}; \emph{Segre} --- l.\ c.

\origpage{265}

The (complete) system $|C'|$, adjoint to the system $|C|$, is bound up with the latter in such a way that if a birational (or Cremona) transformation of the whole plane changes the systems $|C|$ and $|C'|$ into $|D|$ and $|D'|$ (of another plane), \emph{$|D'|$ is again the system adjoint to $|D|$}, provided that there be suitably introduced into $|D'|$ the exceptional curves which come from the transformation${}^{*)}$; in many questions it is not even necessary to attend to this last warning. That is why we may say that \emph{the bond which passes between a linear system of plane curves and its adjoint system is invariable with respect to the birational transformations of the plane.}

It comes from this that the consideration of a system of plane curves together with the successive adjoint systems (namely, the adjoint $|C'|$ which we have just defined, the adjoint of the adjoint, or second adjoint $|C''|$, etc.) may render many services in the study of the Geometry of the plane (or of the rational surfaces), when one takes as fundamental group (in the sense of M.\ \emph{Klein}) the group of the Cremona transformations. It is M.\ S.\ \emph{Kantor} who seems to have perceived first the profit that might be drawn from this method. In fact he showed in a Note of the Comptes Rendus de l'Ac.\ d.\ Sc.\ (1885) how one may approach by this route the problem of the determination of the \emph{cyclic} plane Cremona transformations (transformations of which a suitable power gives the identity); and he has continued this research in other more recent works${}^{**)}$. Later M.\ \emph{Castelnuovo} profited by the adjoint systems to study the linear systems of plane curves${}^{***)}$, and M.\ \emph{Enriques} made use of them for the determination of the continuous groups of birational transformations of the plane${}^{\dagger)}$.

The property of invariance of the adjoint system which we have just recalled permits also the definition of the system $|C'|$ adjoint to a system $|C|$ situated on a rational surface; it will in fact be necessary to choose $|C'|$ in such a way that in any birational representation of the surface on the plane $|C|$ and $|C'|$ should correspond to two systems of plane curves, of which the second is the adjoint to the first.

Now when the rational surface $F$ is given by its projective characters (namely, a surface of ordinary space, of a known order $n$, etc.), one may ask to construct directly on $F$ the system $|C'|$ adjoint

\textbf{*)} \emph{Nöther}, \emph{Rationale Ausführung der Operationen}... §~31, Math.\ Annalen 23, 1883; \emph{Castelnuovo}, \emph{Ricerche generali sopra i sistemi lineari}... n${}^{\text{o}}$ 27.

\textbf{**)} \emph{Premiers fondements pour une théorie}..., 4${}^{\text{me}}$ partie, Atti dell'Accad.\ d.\ Sc.\ fis.\ e mat.\ di Napoli, serie II, vol.\ 4${}^{\text{o}}$. See also «Journal für die Mathem.\ 114», and «Acta Mathem.\ 19».

\textbf{***)} \emph{Ricerche generali sopra i sistemi lineari}..., l.\ c.

\textbf{†)} \emph{Sui gruppi continui}..., Rendic.\ della R.\ Accad.\ d.\ Lincei, 1893.

\origpage{266}
to a given system $|C|$, without having recourse to the plane representation. Suppose, for greater simplicity, that $|C|$ is the system $\infty^{3}$ of the plane sections of $F$; will there exist a system of surfaces cutting out on $F$ the system of curves $|C'|$ adjoint to $|C|$? Certainly; one sees even${}^{*)}$ that this property belongs to certain surfaces $\Phi$ of the order $n-3$, which meet every general plane in the curves of order $n-3$ adjoint to the corresponding plane section of $F$.

It results from this that these adjoint surfaces $\Phi$ are comprised among the surfaces $\Psi$ of the same order sub-adjoint to $F$. Now there is occasion to ask whether the $\Phi$ coincide with the $\Psi$, or whether new conditions must be imposed on the $\Psi$ in order to obtain the $\Phi$.

In certain cases one may effectively verify that the $\Phi$ and the $\Psi$ are the same thing. Thus this happens if the surface $F$ has no \emph{isolated multiple points $P$} or \emph{proper} ones, that is to say points $P$ such that every plane drawn through $P$ cuts on the $F$ a curve having a genus \emph{lower} than the genus of the general plane section of $F$. The general surface of the third order, the surface of the fourth order having a double straight line, or a double conic section, give us examples of such surfaces. The same peculiarity (coincidence of the $\Phi$ with the $\Psi$) may moreover present itself also in other more extended cases${}^{**)}$.

It must not be thought, however, that it is always so. On the contrary, it is proved that the adjoint surfaces have a point of the order $i-2$ at every (ordinary) multiple point of the order $i$ of $F$; while the sub-adjoint surfaces may not pass at all through that point, or else pass through it with a lower order. Whence it results that the system of the adjoint surfaces $\Phi$ may well have a lower dimension than the system of the sub-adjoint surfaces $\Psi$.

Let us verify this remark on an example. Let us consider to this end the surface $F$ of the fourth order having a triple point $P$ (an \emph{isolated} multiple point), and designate by $|C|$ the system of the plane sections of the $F$. One sees easily that the $F$ may be represented on the plane in such a manner that to the system $|C|$ there corresponds a system $|c|$ of plane curves of the fourth order meeting one and the same cubic (the image of $P$) in 12 fixed points. Now the surfaces $\Psi$ of order $4-3=1$ sub-adjoint to $F$ are, according to the definition, the planes of space. These, however, do not cut out on $F$ the system $|C'|$ adjoint to $|C|$; for $|C'|$ is cut out by those planes $\Phi$ alone which pass through the point $P$; one sees it at once if one has recourse to the plane representation, where $|C'|$ is represented by the system $|c'|$ of the straight lines adjoint to the curves $|c|$. It follows therefore that from the

\textbf{*)} \emph{Enriques}, \emph{Ricerche di geometria}... III, 5; \emph{Introduzione}... IV; \emph{Humbert}, \emph{Sur la théorie générale des surfaces unicursales}, Math.\ Annalen 45, 1894.

\textbf{**)} See on this subject the cited Memoirs of MM.\ \emph{Enriques} and \emph{Humbert}.

\origpage{267}
system $\infty^{3}$ of the \emph{sub-adjoint} $\Psi$ one deduces the system $\infty^{2}$ of the \emph{adjoint} $\Phi$, by imposing on the $\Psi$ the condition of passing through the point $P$.

The preceding considerations regard the case of ordinary multiple points. But as soon as higher singularities are in question, the behaviour of the adjoint surfaces can no longer be defined in so simple a manner, although one may always have recourse to the representation of the surface on the plane.

This last means will nevertheless fail when one passes from the rational surfaces to the irrational surfaces. In fact, even for these, one may define a system of \emph{adjoint} surfaces $\Phi$, comprised among the sub-adjoint $\Psi$, and such that the $\Phi$ cut out on the surface a system \emph{adjoint} to the system of the plane sections, enjoying properties analogous to those which we have just seen on the subject of the rational surfaces. The behaviour of the $\Phi$ at an ordinary isolated multiple point of the new surface may be established by the same law which we have set out in the case of the rational surfaces; but the singular multiple points present here a still greater difficulty; for one can no longer have recourse to the representation on the plane.

One sees from this the opportuneness of establishing a new definition of the adjoint system, a definition which should extend to all algebraic surfaces without attending to their projective characters (order, singularities, etc.), and which should permit, in consequence, a deeper study of the bond (invariable with respect to the birational transformations) which passes between a system of curves and its adjoint. To attain this end it is fitting first to examine the properties of the adjoint system in the case where we have altogether defined this system, that is to say on the plane. We shall then seek to extend our conclusions to the systems situated on any surface.

\subsection*{15. On some characteristic properties of the adjoint system.}

Let us take up again, then, a linear system $|C|$ of plane curves of order $m$ and genus $\pi$, and the adjoint system $|C'|$ whose curves have the order $m-3$; and let us try to establish the bond which passes between $|C|$ and $|C'|$ without attending to the projective characters of the two systems.

We remark first that $|C'|$ \emph{cuts out the canonical series} $g_{2\pi-2}^{\pi-1}$ \emph{on every general curve of} $|C|$.

This property does not, however, always suffice to define the system $|C'|$. To see it, one has only to take up again the system $\infty^{3}$ of curves of the fourth order which pass through 12 points of a cubic, a system which represents on the plane the surface of the fourth order having a

\origpage{268}
triple point; in fact the system itself has the property of cutting out on each of its curves the canonical series $g_{4}^{2}$, while the adjoint system $(\infty^{2})$ is formed by the straight lines of the plane.

In spite of this it is fitting for us to consider the widest system $|C_{1}|$ cutting out the canonical series on the general curve of $|C|$; it is proved that this system $|C_{1}|$ is unique and comprises (up to fixed curves) every curve which meets the $C$ in canonical groups. $|C_{1}|$ will not always be the system adjoint to $|C|$; we call it the \emph{system sub-adjoint to} $|C|$, for one sees easily that on a rational surface $F$, whose plane sections are represented by $\infty^{3}$ curves of $|C|$, the system $|C_{1}|$ is cut out by the surfaces sub-adjoint to $F$ (surfaces of the order $M-3$, if $M$ is the order of $F$${}^{*)}$)).

The sub-adjoint system $|C_{1}|$ always contains the adjoint system $|C'|$ (and may even sometimes coincide with it). It is now a question of examining by what conditions one may deduce the adjoint system from the sub-adjoint system. If we take up once again the rational surface $F$ which we have just considered, we see at once that the question proposed does not differ from that which we met above: to establish the conditions which must be imposed on a sub-adjoint surface $\Psi$ in order that it should become an adjoint surface $\Phi$. These conditions depended on the existence on $F$ of certain \emph{isolated} (or \emph{proper}) multiple points, which have no effect on the sub-adjoint surfaces. Now let us remark that an isolated multiple point of $F$ is represented on the plane by a \emph{fundamental curve $\gamma$} of the system $|C|$, a curve which is not met in variable points by the $C$; it is moreover a fundamental curve which we shall call \emph{proper}, having regard to the property that the residual system $|C-\gamma|$ has a genus lower than the genus of $|C|$ (in relation with the fact that the multiple point $P$ is \emph{isolated}). One has therefore to seek on the plane what are the conditions which a proper fundamental curve of the system $|C|$ imposes on the adjoint curves. But this investigation (as well as its analogue on the adjoint surfaces) does not fail sometimes to present serious difficulties.

Here is how one succeeds, by an indirect route, in avoiding them. Let us put our system $|C|$ into relation with any other linear system $|D|$ of curves of the same plane. Let us consider further the system $|C+D|$, the sum of $|C|$ and $|D|$, and the system $|(C+D)'|$ adjoint to it. The general curve of the system $|C+D|$ has the order $m+n$ (if $m$ and $n$ are the orders of $C$ and $D$), and has the order of multiplicity $\mu+\nu$ at every point which is a \emph{base} point of order $\mu$ and $\nu$ for $|C|$ and $|D|$.

\textbf{*)} For example, on the surface of the fourth order having a triple point, the system sub-adjoint to the system of the plane sections is that same system.

\origpage{269}
Consequently the general curve of the system $|(C+D)'|$ has the order $m+n-3$ and the multiplicity $\mu + \nu - 1$ at the point named. Now if we compare this system $|(C + D)'|$ with the system $|C' + D)|$, the sum of $|C'|$ and of $|D|$, we see that the two systems are formed of curves of the same order $(m - 3) + n$, which have the same multiplicity $(\mu - 1) + \nu$ at every point which is a \emph{base} point $(\mu \geqq 1)$ for $|C|$. But a distinction between the two systems presents itself at every point (if such exist) which is a base point for $|D|$, but not for $|C|$ $(\mu = 0, \nu > 0)$; for such a point has the order $\mu + \nu - 1$ for the system $|(C + D)'|$, and the order $\mu + \nu$ for the system $|C' + D|$. According to this we may state the following result:

\emph{If $|C|$ and $|D|$ are two systems of plane curves, the system $|(C + D)'|$ adjoint to their sum contains the system $|C' + D|$; and precisely, if one wishes to deduce $|C' + D|$ from $|(C + D)'|$, one must increase by one unit the order of every point which is a base point for $|D|$, but is not one for $|C|$.}

Now there is occasion to ask: does this property subsist also for the sub-adjoint systems? Not always. And then, may one not make use of the same property to separate the adjoint system from the sub-adjoint system which contains it?

It is precisely so. Let us start from any system $|C|$, and choose an auxiliary system $|D|$ having no proper fundamental curves; (this is possible in an infinity of ways). Now let there be chosen, among the curves of the system $|C_{1}|$ sub-adjoint to $|C|$, those $C'$ which, added to any curve $D$, give curves sub-adjoint to $|C + D|$; it can be proved that these curves $C'$ form a linear system which does not depend on $|D|$; \emph{$|C'|$ is the system adjoint to $|C|$}. One may draw from this remark a new definition of the system adjoint to a given system of plane curves. Now, and this is a very interesting result, this definition may be carried over to the linear systems situated on \emph{any} algebraic surface. One arrives thus at the following statement, of which the reader will find the proof elsewhere${}^{*)}$.

\subsection*{16. System adjoint to a given linear system on an algebraic surface.}

\emph{When one knows on an algebraic surface an irreducible linear system $|C|$, one may construct a second system $|C'|$, called the system adjoint to $|C|$, which is altogether defined by the following properties:}

\textbf{*)} To the definition of the adjoint system here reported is devoted Chap.\ III of the \emph{Introduzione}... of M.\ \emph{Enriques}.

\origpage{270}

1) \emph{the curves of $|C'|$ cut out groups of the canonical series on every general curve of $|C|$;}

2) \emph{whatever the system $|D|$ may be, the system $|C' + D|$ cuts out groups of the canonical series on every general curve of $|C + D|$ (supposed irreducible). These groups however contain $i$ points united at each point which has the order $i$ for $|D|$, and is not a base point for $|C|$.}

It must be well remarked that the series $g_{2\pi-2}$ cut out by the adjoint system $|C'|$ on the general curve $C$ (supposed of genus $\pi$) is not necessarily complete; on some surfaces (contrary to what happens on the plane and on the rational surfaces) the said series may well have a dimension lower than $\pi - 1$. It may even be wanting, together with the adjoint system itself (apart from the evident case $\pi = 0$). Thus it is proved${}^{*)}$ that there is no system adjoint to the linear system of the plane sections of a ruled surface of any genus $\pi$.

It is often fitting to state, under the form of a theorem, the definition which we have just given of the adjoint system, or at least the essential part of the same definition. One arrives thus at the \emph{fundamental property of the adjoint system}:

\emph{Let $|C|$ and $|D|$ be two linear systems of curves on one and the same surface; the system $|C'|$ adjoint to the first, added to the second, gives a system $|C' + D|$ which is contained in the system $|(C + D)'|$ adjoint to $|C + D|$; it differs from it only in this, that there may exist points at each of which the order of multiplicity of $|C' + D|$ is higher than the corresponding order of $|(C + D)'|$.} These are precisely the points which are base points, of order $\nu$ for example, for $|D|$, while they are not base points for $|C|$; in fact at each of these points $|C' + D|$ has the order of multiplicity $\nu$, while $|(C + D)'|$ has the order of multiplicity $\nu - 1$. If in particular these points are wanting, so that every base point of $|D|$ is also a base point of $|C|$, one has
\[
|C' + D| = |(C + D)'|.
\]

According to the preceding theorem one may, by a simple addition, construct the system $|(C + D)'|$ adjoint to $|C + D|$, as soon as one knows the system $|C'|$ adjoint to $|C|$. Inversely, to deduce the system adjoint to $|C|$ from the system adjoint to $|C + D|$, one has only to subtract the system $|D|$ from the last adjoint system. From these observations it easily results that, as soon as one knows the system $|C'|$ adjoint to a particular system $|C|$ on a surface, one may construct the system adjoint to a system arbitrarily given on the surface, by the aid of simple operations of addition and of subtraction (with impositions or suppressions of multiple points, if it is necessary).

\textbf{*)} \emph{Enriques}, \emph{Introduzione}..., 30.

\origpage{271}

\emph{Remark.} --- The preceding definition, as has been seen, does not establish in a direct manner the characters of the operation of \emph{adjunction}, by which one passes from a given system $|C|$ to the adjoint system $|C'|$. On the contrary this operation is defined by its behaviour in relation with another operation known beforehand (addition). This behaviour is represented by the relation
\[
|(C + D)'| = |C' + D|
\]
which always subsists, up to multiplicities at certain points. Now in Analysis one often has occasion to have recourse to such definitions; every time that one defines a function, not directly, but by the aid of a \emph{functional equation}. This comparison with known theories may serve to bring to light the route which we have just followed.

\subsection*{17. Direct definition of the adjoint system.}

In the simplest cases one may give a direct definition of the system adjoint to a given system, a definition which follows from the preceding one. Let us confine ourselves to the following statement:

\emph{Let $|C|$ be a system which possesses no proper fundamental curves} (that is to say curves $\gamma$ which are not met in variable points by the $C$, and such that $|C - \gamma|$ has a genus lower than the genus of $|C|$); \emph{then, in general, every curve which cuts out a group of the canonical series on the curves $C$ belongs to the system adjoint to $|C|$.} The words \emph{in general} relate to the following observations.

To prove the theorem one supposes first that the curves $C$ which pass through any point of the surface do not in consequence pass through other points mobile with that one. One meets next a second restriction; it relates to the case where, the surface belonging to the \emph{class} of the ruled surfaces, the curves $C$ would cut the generators in one point. While the first restriction is perhaps superfluous, the second is certainly necessary. In fact the system $|C|$ of the plane sections of a ruled surface has no adjoint system, but there are indeed curves which cut out on the $C$ groups of the canonical series.

Returning to the direct definition of the adjoint system, there would now be the case of a system $|C|$ endowed with proper fundamental curves to consider. In order not to weary the reader we shall not insist on this subject; for these developments one may consult n${}^{\text{o}}$ 29 of the \emph{Introduzione}... of M.\ \emph{Enriques}.

\origpage{272}

\subsection*{18. Surfaces adjoint to a given surface.}

The definition of the system adjoint to a given system (n${}^{\text{o}}$ 16) has just been stated under a form which undergoes no change; when one transforms birationally the surface considered; a system $|C|$ and its adjoint $|C'|$ are transformed into two systems $|D|$ and $|D'|$ of which the second is again the adjoint to the first (up to \emph{exceptional} curves). That is why this definition lends itself better than any other in researches into the properties of invariance of an algebraic surface; we shall see it presently.

But first it is fitting for us to make some remarks from the projective point of view.

Let $|C|$ be the system of the plane sections of a surface $F$ of order $n$ of ordinary space; and let $|C'|$ be the system adjoint to $|C|$. Then one may construct a linear system of surfaces $\Phi$ of order $n - 3$, sub-adjoint to $F$, which cut out on $F$ (outside the multiple curves of $F$) the curves of $|C'|$. These surfaces $\Phi$ are altogether determined by the curves $C'$; they are called \emph{surfaces adjoint to $F$}, of the order $n - 3$; we have already spoken of them in the case where $F$ was a rational surface, and we said then that this definition could be extended to general surfaces.

Without insisting anew on the same subject, let us confine ourselves to remarking that the preceding considerations furnish an indirect definition of the surfaces $\Phi$, by the aid of the system $|C'|$ which these cut out on $F$. It results from the definition that the $\Phi$ are comprised among the sub-adjoint surfaces $\Psi$ of the same order $n - 3$, and that they differ from them only by their behaviour at the isolated multiple points of the $F$. So that if the $F$ possesses no such points, the adjoint surfaces coincide with the sub-adjoint; they are then altogether determined by the condition of cutting out on every general plane curves adjoint to the corresponding section of $F$. This is what happens for example if the $F$ is a general surface of the order $n$, or if it possesses a double curve and a finite number of points of order $k = 3, 4, \ldots$ satisfying the condition of having the order $\dfrac{k(k-1)}{2}$ for the double curve.

If on the contrary the $F$ possesses \emph{ordinary isolated} multiple points of the order $k$, it is proved that the $\Phi$ pass with the order $k - 2$ through each of them. So that one may state the following result:

\emph{If a surface $F$ has no other singularities than ordinary multiple curves of the order $h = 2, 3 \ldots$, and ordinary multiple points of the order $k = 2, 3 \ldots$, the adjoint surfaces are altogether determined by the conditions of passing $h - 1$ times through these curves and $k - 2$ times}

\origpage{273}
\emph{through these points.} In particular a double, isolated, ordinary point presents no condition to the adjoint surfaces.

The preceding statement gives us a direct definition of the surfaces adjoint to a surface endowed with ordinary singularities. One might seek to extend this definition to surfaces having singular multiple points; but one meets some difficulties if one wishes to comprise, under a single statement, all the cases which these points may present. That is why we have followed in this connection an indirect route, in defining the surfaces $\Phi$ adjoint to $F$ by the aid of the system $|C'|$ which they cut out on $F$. But one might have approached the problem of the behaviour of the surfaces $\Phi$ at a singular point by leaning on analytic considerations, as M.\ \emph{Nöther} has shown${}^{*)}$. This last route (which from the historical point of view precedes ours) would then permit the adjoint system $|C'|$ to be defined by the aid of the $\Phi$, while we have followed the inverse path.

\emph{Remark.} --- Just as we have done on the plane, so one might define on any surface the curves $C_{1}$ \emph{sub-adjoint} to a system $|C|$, by the single condition of cutting out on every curve $C$ groups of the canonical series. These curves form a linear system $|C_{1}|$ (\emph{sub-adjoint} to $|C|$), with a single exception relative to the ruled surfaces. The system $|C_{1}|$ contains the adjoint system $|C'|$.

If $|C|$ is the system of the plane sections of a surface $F$ of the order $n$ of ordinary space, it is proved that $|C_{1}|$ is cut out by the surfaces of order $n-3$ sub-adjoint to $F$; the case excepted where $F$ is a ruled surface of genus $\pi > 0$, for then there are $C_{1}$ composed of $2\pi-2$ generators, but there exists no surface of the order $n-3$ sub-adjoint to $F$ which meets $F$ in these generators only (outside the multiple curves).

For the rest, we shall not have to consider, in what follows, the sub-adjoint system.

\subsection*{19. Adjoint surfaces of other orders.}

When one has defined, by the means which have just been seen, the surfaces $\Phi$ of order $n-3$ adjoint to $F$, one may also construct \emph{adjoint surfaces of order} $n-3+i$ $(i \gtreqqless 0)$, that is to say surfaces which have

\textbf{*)} \emph{Zur Theorie}... Math.\ Annalen 8, pag.\ 510; see also a Note of the same author in the Götting.\ Nachrichten (1871), and the Memoir \emph{Ueber die totalen algebraischen Differentialausdrücke}, Math.\ Annalen 29. M.\ \emph{Nöther} is occupied in particular with the adjoint surfaces of the order $n-4$ ($\varphi$-\emph{Flächen}), which we shall study next. See also (for particular surfaces) the cited Memoir of M.\ \emph{Humbert}, and the other Memoir of the same author \emph{Théorie générale des surfaces hyperelliptiques}, Journal de Mathém.\ 4${}^{\text{e}}$ s${}^{\text{e}}$, t${}^{\text{e}}$ IX, pag.\ 394, 1893.

\origpage{274}
this order, and which behave like the $\Phi$ along the multiple curves and at the multiple points of $F$. The new surfaces are also determined by the condition of cutting out on $F$ the curves of the system $|C'+iC|$, where $|C|$ is the system of the plane sections of $F$, and $|C'|$ is the system adjoint to it. Now according to the fundamental theorem on the adjoint system, $|C'+iC|$ is nothing else than the system adjoint to the system $|(i+1)C|$, the multiple of $|C|$. One has in this remark a means of defining the system $|C'+iC|$ (and consequently the adjoint surfaces of order $n-3+i$ which cut it out) every time that $|C'|$ is wanting; for even then the system $|(i+1)C|$ may well have an adjoint system; and it certainly has one (as is proved) when $i$ exceeds a certain value.

We shall meet in what follows these systems $|C'+iC|$; we shall then have to have recourse to the remark which we make from now on. We have already remarked (n${}^{\text{o}}$ 17) that the system $|C'|$ adjoint to a system $|C|$ of genus $\pi$ cuts out on the general curve $C$ a series $g_{2\pi-2}$, which is contained in the canonical series $g_{2\pi-2}^{\pi-1}$, but which may well be incomplete and have consequently a certain \emph{defect}; it may even be lacking (the example of the ruled surfaces has proved it to us), and in that case the defect attains its \emph{maximum} $\pi$. Likewise the series which the system $|C'+iC|$ $(i > 0)$ cuts out on $|C|$ is a series $g_{2\pi-2+in}$ (where $n$ designates the degree of $|C|$); if it were complete, its dimension would be $\pi-2+in$; but it may happen that it too is incomplete, and then one will have to consider its defect. Now, and this is the remark which we proposed to make, the consideration of this defect takes place only for the first values of $i$; for \emph{one may always fix such a limit $l$ that for the values of $i > l$ the series cut out by $|C'+iC|$ on $C$ results complete.}

\subsection*{20. Surfaces adjoint to a skew curve.}

In n${}^{\text{o}}$ 18 we have learnt to construct surfaces which cut out on a given surface $F$ of order $n$ the system adjoint to the plane sections of the latter. Likewise, by the aid of the fundamental theorem, one may construct surfaces which cut out on $F$ the system adjoint to \emph{any} system of curves $|C|$ given on $F$. To this end one first makes pass through the general curve $C$ a surface $G$ of an order $m$ high enough for this construction to be possible. Let us call $K$ the curve residual intersection of the surface $G$ with $F$; one will then have $|K+C| = |mC|$. Now the system $|C'+(m-1)C|$ adjoint to it is cut out on $F$ by the surfaces of order

\origpage{275}
\[ n-3+(m-1) = m+n-4 \]
adjoint to $F$ (n${}^{\text{o}}$ 19). Therefore, according to the fundamental theorem:

\emph{If we impose on these adjoint surfaces of order $m+n-4$ the condition of passing through the curve $K$, and of having with $F$ a contact of order $i-1$ at every base point of order $i$ for the system $|C|$, we shall cut out on $F$ the system adjoint to $|C|$.}

The theorem which has just been stated furnishes us first of all the means of cutting out on a skew curve $C$, partial intersection of two surfaces $F$ and $G$ having the orders $n$, $m$, the canonical series, or at least a part of that series. We see that for this purpose one must construct the surfaces of order $m+n-4$ which pass through the residual intersection $K$ of $F$ and $G$, and behave in a suitable manner at the multiple points of these surfaces. Now this is precisely a theorem of M.\ \emph{Nöther} on the surfaces \emph{adjoint} to a skew curve${}^{*)}$.

But it results further from the preceding theorem that every curve adjoint to $|C|$ on the surface $F$ may be cut out by the said adjoint surfaces, which permits, so to say, the inversion of the theorem of M.\ \emph{Nöther}, so far as it is possible.

\section*{Chapter V. Invariants of a surface with respect to the birational transformations.}

Leaning on the preceding results we are now in a position to approach the fundamental question of our theory; namely, the search for some invariants of a surface with respect to the birational transformations; invariants which may be numbers or geometric varieties (curves or systems of curves).

It is fitting, however, to let this search be preceded by some remarks of a general character.

The notion of the invariants was arrived at first of all by considering the transcendental expressions bound up with the surface (\emph{Clebsch}, \emph{Nöther} 1868---69; \emph{Picard} 1884). Next it was proposed to define the same invariants by the geometric (or algebraic) route. The geometers who approached this question (\emph{Cayley}, \emph{Nöther}, \emph{Zeuthen}) proceeded in the following manner. A surface $F$ of ordinary space was considered, and attention was paid to its projective characters (order, multiple curves and points). One then transformed birationally

\textbf{*)} \emph{Zur Theorie}... pag.\ 510, Math.\ Annalen 8; see also the Memoir, \emph{Ueber die algebraischen Raumcurven}, Denkschriften der Berliner Akad.\ 1882, pag.\ 12.

\origpage{276}
the $F$ into another surface $F'$ of the same space, and by comparing the projective characters of $F$ with the analogous characters of $F'$ one formed numerical expressions or functions which were not modified by the transformation.

We shall now replace this method by another which presents some advantage over it, although it does not differ from it in an essential manner; it permits us nevertheless to abstract from the projective characters of the surfaces, and thereby to avoid the difficulties which come from the singular points and from their behaviour in the transformations.

To this end let us take up again the general surface of a given \emph{class} (n${}^{\text{o}}$ 3), a surface whose projective characters have for us no interest, and let us start from a linear system of curves $|C|$ of the surface. From $|C|$ we may deduce, by some operation, new systems which are bound up with $|C|$; such are, for example, the adjoint system $|C'|$, the adjoint $|C''|$ of $|C'|$, etc. Let one now seek relations between these systems, or between their characters (genus, dimension, degree); if one arrives at a relation which does not depend on the system $|C|$ from which one starts, one will have there an \emph{invariant of the surface}.

The examples which we shall give will show the fruitfulness of this method.

\subsection*{21. Canonical system.}

The most important example is furnished us by the immediate application of the fundamental theorem on the adjoint system.

This theorem tells us in fact that if $|C|$ and $|D|$ are two systems of curves on the surface, one arrives at the system $|(C+D)'|$ adjoint to $|C+D|$ in two different manners: one may add to $|D|$ the system $|C'|$ adjoint to $|C|$, but one may likewise add to $|C|$ the system $|D'|$ adjoint to $|D|$. One has therefore the symbolic equality${}^{*)}$
\[ |C'+D| = |C+D'|, \]
whence one deduces
\[ |C'-C| = |D'-D|. \]
This relation has a geometric signification only in the case where $|C'|$ contains $|C|$, for it tells us then that $|D'|$ also contains $|D|$, and moreover that the system $|K| = |C'-C|$ residual to $|C|$ with

\textbf{*)} On the subject of the orders of multiplicity of the base points which belong to the two systems considered in this relation there would be an observation to make, for these orders may differ by one unit; let us be permitted to omit it, seeing that here we propose only to sketch the proof. See the complete proof in n${}^{\text{o}}$ 38 of the \emph{Introduzione}... of M.\ \emph{Enriques}.

\origpage{277}
respect to $|C'|$ is likewise the residual to every other system $|D|$ with respect to its adjoint $|D'|$. In conclusion:

\emph{The residual $|K|$ of a linear system $|C|$ with respect to its adjoint system does not depend on the system $|C|$ from which one starts.}

In other terms, $|K|$ is a system of curves which depends only on the surface considered, or more precisely on the \emph{class} of surfaces to which the latter belongs; it is therefore an \emph{invariant} of the said surface. We shall call it the \emph{canonical system} (and canonical curve each of its curves), for, as we are going to see, it has the greatest analogy with the canonical series defined on curves.

The dimension of the canonical system $|K|$ is also an \emph{invariable character} of the surface. This dimension increased by one unit is called the \emph{first genus} of the surface (\emph{Flächengeschlecht} after M.\ \emph{Nöther})${}^{*)}$; it is fitting to add \emph{geometric genus}, to distinguish it from the numerical genus which we shall soon introduce; it is designated by $p_{g}$.

If $p_{g} = 0$, there is no canonical system; in other terms, there is on the surface no system $|C|$ contained in its adjoint $|C'|$. This is what happens, for example, on the rational surfaces, on the ruled surfaces, or more generally on every surface containing a system $|C|$ of curves whose characteristic series is non-special.

If $p_{g} = 1$, the general curve of any system $|C|$ belongs to \emph{a single} curve of the adjoint system $|C'|$. The remainder $C' - C$, if there be any, furnishes (up to exceptional curves) the single canonical curve. But in certain cases it may happen that $|C|$ coincides with $|C'|$; then there is no remainder, and the canonical curve has, so to say, the order zero.

If $p_{g} \geqq 1$, the canonical system $|K|$ may be constructed by the aid of a single system of curves $|C|$ of the surface, and of the adjoint system $|C'|$; this results from the definition. In recalling the properties of the adjoint system, one finds for the canonical curves the following property:

\emph{Every canonical curve cuts out on the general curve $C$ of any system $|C|$ a group of points which, joined to the base points of $|C|$ and to a group of the characteristic series of $|C|$, furnishes a group of the canonical series of $C$.}

The property stated is characteristic for the canonical curves; it even suffices, in order to define these curves, that it be verified with respect to a single system $|C|$, provided that this system satisfy the restrictions spoken of in n${}^{\text{o}}$ 17.

\textbf{*)} \emph{Zur Theorie}... pag.\ 520, ibid.

\origpage{278}

\subsection*{22. Construction of the canonical system by the aid of the adjoint surfaces of the order $n-4$.}

When a surface $F$ is given by its projective characters in ordinary space, and one designates by $n$ its order, it is fitting, in the construction of the canonical system $|K|$, to have recourse to the system $|C|$ of the plane sections of $F$. One has to subtract the system $|C|$ from the adjoint system $|C'|$, which is cut out on $F$ by the adjoint surfaces of order $n-3$. The residual system will naturally be determined by the adjoint surfaces of order $n-4$. One has therefore this result:

\emph{On a surface of order $n$ of ordinary space the canonical system is cut out by the adjoint surfaces of the order $n-4$}${}^{*)}$. The number of such surfaces which are linearly independent therefore furnishes the genus $p_{g}$ of $F$. Here, as ordinarily, one considers the intersection of an adjoint surface with $F$, outside the multiple curves of $F$. But in this intersection one may sometimes find curves which belong to all the adjoint surfaces of order $n-4$, and which are not to be regarded as canonical curves. It is proved in fact that if a curve $\gamma$ of $F$ is \emph{exceptional}, that is to say if it corresponds, by a birational transformation, to a point $P'$ of a surface $F'$ (n${}^{\text{o}}$ 2), this curve $\gamma$ belongs to all the surfaces of order $n-4$ adjoint to $F$.

Now on the subject of the exceptional curves $\gamma$ of $F$ there is a distinction to be made. In fact the point $P'$ of $F'$ which corresponds to the curve $\gamma$ does not in general belong to all the canonical curves of $F'$ (is not a base point of the canonical system on $F'$), and in this case the exceptional curve $\gamma$, which we shall call of the \emph{first kind}, is not to be regarded as forming part of the canonical curves of $F$; consequently one must subtract the $\gamma$ from the intersection of $F$ with a general adjoint surface of the order $n-4$ in order to obtain a canonical curve. It will be further remarked that an exceptional curve of the first kind has no variable intersection with the canonical curves${}^{**)}$.

But it may on the contrary happen, in some particular cases, that $P'$ is a base point of the canonical system on $F'$; then the corresponding exceptional curve $\gamma$ on $F$, an exceptional curve of the \emph{second kind}, enjoys the property of invariance, and is to be regarded as forming part of all the canonical curves. The variable residual

\textbf{*)} \emph{Nöther}, \emph{Zur Theorie}... pag.\ 514, ibid.

\textbf{**)} To have an example of such an exceptional curve it suffices to consider the straight line which joins two triple points of a surface of the fifth order; here the adjoint surfaces of the order $5-4=1$ are the planes drawn through this straight line.

\origpage{279}
curves now cut the $\gamma$ in variable points (one at least).

Here is how one may obtain an example of an exceptional curve of the second kind. Let one consider a surface of the fifth order having two contacts with itself (\emph{tacnodal points})${}^{*)}$; each of these points imposes a condition on the adjoint surfaces, so that the adjoint surfaces of the order $5-4=1$ are the planes passing through these points. The canonical curves are therefore curves of the fifth order (genus 2), which have a base point $P'$ (outside the singular points) in the further intersection of the surface with the straight line joining the two singular points. Every birational transformation of the surface which changes the point $P'$ into a curve $\gamma$ gives rise to an exceptional curve of the second kind.

Leaving aside the example to return to the preceding remark, one may now make more precise the statement of the last theorem:

\emph{The surfaces of order $n-4$ adjoint to a surface of the order $n$ cut out on the latter the canonical system, outside the multiple curves and the exceptional curves of the first kind.}

The consideration of the adjoint surfaces of the order $n-4$ might even serve to define the canonical system. It is precisely the route which \emph{Clebsch} indicated, and which M.\ \emph{Nöther} has followed in the cited Memoir. This savant had naturally to surmount the difficulty of proving the invariant character of the system which he thus came to construct. We have avoided this difficulty by defining the canonical system by the aid of the fundamental property of the adjoint system, which may even replace the property of invariance of the canonical system for the surfaces $(p_{g} = 0)$ which do not possess such a system.

The construction of the canonical system by the aid of the adjoint surfaces of the order $n-4$ shows that this system has the greatest analogy with the canonical series cut out on a plane curve of the order $n$ by the adjoint curves of the order $n-3$; one sees thus the reason of the name which we have attributed to this system.

\subsection*{23. Reducible canonical system.}

The canonical system belonging to a given surface may be reducible, either by the presence of some fixed curve common to

\textbf{*)} If 1 and 2 (fundamental points of the system of coordinates) are the two singular points, the equation of the surface is of the form
\[
a x_{1}^{3} x_{2}^{2} + b x_{2}^{3} x_{1}^{2} + x_{1}^{2} x_{2}^{2} f_{1} + x_{1} x_{2} (x_{1} g_{2} + x_{2} h_{2} + i_{3}) + x_{1} j_{4} + x_{2} k_{4} + l_{5} = 0
\]
where $a$ and $b$ are constants, and $f_{1}, g_{2} \ldots l_{5}$ are algebraic forms in the variables $x_{3}$, $x_{4}$ having the degrees $1, 2 \ldots 5$.

\origpage{280}
every canonical curve, or by the division of the general canonical curve into several curves variable in a pencil (n${}^{\text{o}}$ 7).

The first case presents itself, as we know, every time that there is on the surface some exceptional curve of the second kind. But there is occasion to investigate whether the canonical system may not contain other fixed curves besides these exceptional curves. It appears that it may in fact contain some, although it is not easy to find examples on this subject. Without stopping on the question, we confine ourselves to showing, by an example, that for the surfaces having $p_{g} = 1$ the canonical curve is not necessarily exceptional, as one has been led to believe; on the contrary it may happen that this curve (or some one of its parts) has the genus $> 0$. Let us consider in fact a surface of the fifth order passing through a conic section, and having three contacts with itself at three points $A$, $B$, $C$ of this curve. The unique adjoint surface of the order $5-4=1$ is here the plane $ABC$, which meets the surface along the conic section named and along a further curve of the third order passing simply through $A$, $B$, $C$. Now it is clear that this last curve (having the genus $> 0$) is certainly not exceptional.

It may further happen, we have said, that the general canonical curve (apart perhaps from fixed parts) decomposes into several curves (at least $p_{g} - 1$) variable in one and the same pencil (which may be rational or else irrational). If the pencil has no base points, and if there exists no exceptional curve of the second kind, it is proved (after M.\ \emph{Nöther}${}^{*)}$) that \emph{each of the partial curves has the genus} 1. The restriction stated is necessary, for there are surfaces of genus $p_{g} = 2$ having a pencil of canonical curves of genus $> 1$ (with base points). One has only to recall the example of the surface of the fifth order cited in n${}^{\text{o}}$ 22. Inversely \emph{if a surface $(p_{g} \geqq 2)$ contains a pencil of curves of the genus} 1, \emph{every canonical curve decomposes into $i \geqq p_{g} - 1$ curves of this pencil.}

\subsection*{24. The characters $p^{(1)}$ and $p^{(2)}$ of a surface.}

To obtain numbers which enjoy the property of invariance with respect to a given surface, one has only to consider the other characters of the canonical system, beside the dimension $p_{g} - 1$ of which we have already spoken. Thus present themselves the genus and the degree of the canonical system.

\textbf{*)} \emph{Zur Theorie}... pag.\ 523, ib.\ See also \emph{Castelnuovo}, \emph{Osservazioni intorno alla geometria sopra una superficie}, Rendic.\ Istituto Lombardo 1891. Nota I, §~2.

\origpage{281}

The genus (of the general canonical curve), on the supposition $p_{g} \geqq 1$, furnishes us the invariant which is called the \emph{second genus} (\emph{Curvengeschlecht} after M.\ \emph{Nöther}${}^{*)}$) and which is designated by $p^{(1)}$. It may be calculated directly for those surfaces alone which have true canonical curves (of order $> 0$)${}^{**)}$; for the other surfaces it could not be defined save by the aid of a convention; this is what we shall do presently, at least in certain cases.

But first let us speak also of the \emph{third genus} $p^{(2)}$${}^{***)}$, which is the degree of the canonical system, that is to say the number of variable intersections of two canonical curves; in order that this definition should have a sense by itself, one must suppose that one has $p_{g} > 1$ (for $p_{g} = 2$ one will have $p^{(2)} = 0$), and moreover that the general canonical curve be irreducible.

To this last character, however, one does not ordinarily attend, for M.\ \emph{Nöther} has shown${}^{\dagger)}$ that it is bound to $p^{(1)}$ by the very simple relation
\[
p^{(2)} = p^{(1)} - 1.
\]

But on the subject of this relation a remark is necessary; for it is not true without restriction, if by $p^{(1)}$ and $p^{(2)}$ one understands the genus of the canonical curve (supposed irreducible), and the number of \emph{variable} intersections of two canonical curves. This is what the following examples will show us.

Let us take up again to this end the surface of the fifth order which has two points of contact with itself (a surface of which we have spoken in n${}^{\text{o}}$ 22). The canonical curves are here the sections of the surface determined by the planes passing through the singular points; they are consequently $\infty^{1}$ curves of the genus $p^{(1)} = 2$; and yet one has $p^{(2)} = 0$, for two curves of this pencil have no \emph{variable} intersection. The preceding relation $p^{(2)} = p^{(1)} - 1$ is therefore at fault.

Another example, relative to a higher value of $p_{g}$, will not be useless. Let us designate by $1, 2, \ldots, 8$ the eight intersections of three surfaces of the second order, and construct a surface $F^{6}$ of order 6 having a multiple point of order three at $1, 2, \ldots, 7$ and a simple point at 8; this is always possible. Now one sees that the surfaces of order $6 - 4 = 2$ adjoint to $F^{6}$ are the surfaces of the second order which pass through $1, 2, \ldots, 7$, and consequently through 8. Here one has $p_{g} = 3$, $p^{(1)} = 4$, $p^{(2)} = 2 = p^{(1)} - 2$.

\textbf{*)} \emph{Zur Theorie}... pag.\ 520, ibid.

\textbf{**)} Thus for example for the general surface of order 5 one has
\[
p_{g} = 4, \quad p^{(1)} = 6.
\]

\textbf{***)} \emph{Nöther} l.\ c.

\textbf{†)} l.\ c.

\origpage{282}

One sees by these examples that the relation $p^{(2)} = p^{(1)} - 1$ is at fault, and must be replaced by the inequality $p^{(2)} < p^{(1)} - 1$, every time that the canonical system possesses on the surface base points outside the multiple points, or, what comes to the same, if it contains exceptional curves of the second kind.

\subsection*{25. Numerical definition of $p^{(1)}$ and $p^{(2)}$.}

The invariant numbers of a surface of which we have spoken up to now, namely $p_{g}$, $p^{(1)}$ and $p^{(2)}$, have been introduced by the direct consideration of the canonical system and of its geometric characters. Now beside these \emph{geometric} invariants it is fitting, in the theory of surfaces, to consider other invariants defined by numerical expressions formed with the characters of any system of curves on the surface. These new invariants, which we shall call \emph{numerical} or \emph{virtual}, are ordinarily such that they coincide with the corresponding geometric invariants when certain conditions of \emph{regularity} of the surface are fulfilled. They furnish on the contrary new characters of invariance for the \emph{irregular} surfaces, and sometimes they may replace the geometric invariants in the cases where the definition of the latter no longer has a sense. The most important of these numerical invariants is the (first) \emph{numerical genus} $p_{n}$, which corresponds to $p_{g}$; we shall devote to it some of the following paragraphs. Let us occupy ourselves first with the numerical characters $p_{n}^{(1)}$, $p_{n}^{(2)}$, which present themselves beside $p^{(1)}$ and $p^{(2)}$.

Let us start from the hypothesis that the canonical system $|K|$ of our surface exists $(p_{g} > 0)$ and is irreducible; and let us consider on the surface any irreducible system $|C|$, as well as its adjoint system $|C'|$. For greater simplicity let us suppose further that $|C|$ and $|C'|$ on our surface have no base points, and that the surface possesses no exceptional curve: one will then have exactly
\[
|K| = |C' - C|.
\]
One may calculate the genus $p_{n}^{(1)}$ and the degree $p_{n}^{(2)}$ of $|K|$ as a function of the analogous characters $\pi$, $n$ of $|C|$ and $\pi'$, $n'$ of $|C'|$, provided one has recourse to the relations of n${}^{\text{o}}$ 9. Thus one will easily find
\[
\left\{
\begin{aligned}
p_{n}^{(1)} &= \pi' + n - 2\pi + 3,\\
p_{n}^{(2)} &= n' + n - 4\pi + 4.
\end{aligned}
\right.
\tag{1}
\]

On the hypotheses in which we have placed ourselves, the two characters $p_{n}^{(1)}$ and $p_{n}^{(2)}$ do not differ from $p^{(1)}$ and $p^{(2)}$, and consequently enjoy the property of invariance. But it must be remarked that the preceding relations

\origpage{283}
keep a determinate sense even when one abstracts from the restrictions imposed on the canonical system $|K|$. This system may well be wanting $(p_{g} = 0)$, or, though existing, it may differ from the system $|C' - C|$ by the presence of exceptional curves, it may even be reducible; nevertheless one will always be able to construct the expressions (1) by the aid of a system $|C|$ and of its adjoint $|C'|$. But one can no longer affirm now that the values (1) coincide with the values of $p^{(1)}$ and $p^{(2)}$, which may even have no sense at all. There is therefore to examine whether, even in this case, the expressions (1) keep the property of invariance. It is indeed so. One proves in fact the following theorem${}^{*)}$:

\emph{Let $F$ be a surface possessing no exceptional curve, and $|C|$ any system having no base points on $F$; the values $p_{n}^{(1)}$, $p_{n}^{(2)}$ calculated by the aid of the characters of the system $|C|$ do not change if one replaces this system by another satisfying the same condition;} $p_{n}^{(1)}$, $p_{n}^{(2)}$ are consequently two \emph{invariants} of the surface${}^{**)}$. \emph{They always satisfy the relation}
\[
p_{n}^{(2)} = p_{n}^{(1)} - 1.
\]

Let one have on the contrary a surface $F$ possessing exceptional curves; one will likewise be able to define the characters $p_{n}^{(1)}$, $p_{n}^{(2)}$ of $F$, provided one succeeds in transforming birationally the surface $F$ into another $F'$ which has no such curves; for it will then suffice to consider, for example, on $F'$ the system of the plane sections. Now this transformation of $F$ into $F'$ is possible in many cases, even if the canonical system is wanting $(p_{g} = 0)$. Not always, however; there are in fact surfaces (for example, the rational surfaces, and the ruled surfaces) which certainly possess exceptional curves, in whatever manner one transforms them. For these surfaces the definitions of $p_{n}^{(1)}$ and $p_{n}^{(2)}$ which we have given are at fault. One might indeed modify them so that they should extend to these cases as well; but it is not fitting for us to stop here on researches which are still incomplete.

\textbf{*)} \emph{Enriques}, \emph{Introduzione}... n${}^{\text{o}}$ 41.

\textbf{**)} On the surface of the fifth order having two contacts with itself (n${}^{\text{o}}$ 22, 24) a system such as $|C|$ is furnished by the plane sections of the surface $(\pi = 6, n = 5; \pi' = 12, n' = 16)$; one will find $p_{n}^{(1)} = 2$, $p_{n}^{(2)} = 1$, while one had $p^{(1)} = 2$, $p^{(2)} = 0$.

Another example worthy of remark is given us by the surface of the sixth order which passes twice through the six edges of a tetrahedron. This surface, of which we shall have to speak later, has no canonical system $(p_{g} = 0)$ and consequently $p^{(1)}$ and $p^{(2)}$ have no sense; yet if one starts from the system of the plane sections of the surface $(\pi = 4, n = 6; \pi' = 4, n' = 6)$ one finds $p_{n}^{(1)} = 1$, $p_{n}^{(2)} = 0$.

\origpage{284}

\subsection*{26. The numerical genus.}

We shall now define a new numerical invariant of a surface, which also enjoys a great importance.

Let us start from a system $|C|$ situated on the surface, and from its adjoint system $|C'|$. We have remarked formerly (n${}^{\text{o}}$ 16) that the series $g_{2\pi-2}$ cut out by $|C'|$ on the general curve $C$ of $|C|$ is not necessarily complete; on the contrary it may have a certain defect $\delta_{0}$, and then its dimension reduces to $\pi - 1 - \delta_{0}$ (and the dimension of $|C'|$ reduces to $p_{g} + \pi - 1 - \delta_{0}$).

Likewise the series cut out by the systems $|C' + C|$, $|C' + 2C| \ldots$ on the curve may have positive defects $\delta_{1}$, $\delta_{2}, \ldots$ The succession of the positive numbers $\delta_{0}$, $\delta_{1}$, $\delta_{2}, \ldots$ is not, however, unlimited, for we know (n${}^{\text{o}}$ 19) that one may always choose a number $i$ high enough for the systems $|C' + iC|$, $|C' + (i+1)C| \ldots$ to cut out on $C$ a complete series. There is therefore occasion to consider the sum
\[
\Delta = \delta_{0} + \delta_{1} + \cdots + \delta_{i-1}
\]
formed with the defects of the series cut out on the general curve of a system $|C|$ by the systems $|C'|$, $|C' + C| \ldots$ which are adjoint to $|C|$, $|2C| \ldots$.

Well then, this sum $\Delta$, which is altogether determined by the system $|C|$, does not change when one replaces the system $|C|$ by any other system of the surface; in other terms, $\Delta$ is a \emph{new invariant} of the surface${}^{*)}$. Although $\Delta$ itself enters into several questions, nevertheless, conforming ourselves to established usage, we shall introduce in place of $\Delta$ a new character, the \emph{numerical genus} $p_{n}$, which is defined by the relation
\[
p_{g} - p_{n} = \Delta.
\]

\emph{One obtains therefore the numerical genus of a surface by subtracting from the geometric genus the sum of the defects of the series which are cut out on the general curve of a linear system by the systems adjoint to it and to its multiples.}

According to this definition one has
\[
p_{n} \leqq p_{g};
\]
$p_{n}$ may even be negative, while $p_{g}$ never is.

For example a ruled surface having its plane sections of genus $\pi$ has, as we have remarked, $p_{g} = 0$; moreover, if $|C|$ is the system of the plane sections, one finds $\delta_{0} = \pi$, $\delta_{1} = \delta_{2} = \cdots = 0$, and

\textbf{*)} \emph{Enriques}, \emph{Introduzione}... n${}^{\text{o}}$ 40.

\origpage{285}
consequently $\Delta = \pi$, $p_{n} = -\pi$; \emph{the numerical genus of a ruled surface is the genus of the plane sections of the latter preceded by the negative sign.}

The definition of $p_{n}$ which we have just given may also be presented in the projective form; we shall see thereby that this new character is not distinguished from the numerical genus introduced by \emph{Cayley}, \emph{Zeuthen} and \emph{Nöther}${}^{*)}$, and this will give the reason of the name which we have attributed to $p_{n}$.

It suffices, to this end, to suppose that the system $|C|$ of which we make use to calculate $\Delta$ is formed by the plane sections of a surface $F$ of ordinary space, whose projective characters, the order $n$, etc., we now consider. The systems $|C'|, |C'+C|, \ldots, |C'+iC|$ are then cut out on $F$ by the surfaces $\Phi^{\nu}$ of order $\nu = n-3$, $n-2, \ldots n-3+i$ adjoint to $F$. The remark according to which the systems $|C'+iC|, |C'+(i+1)C| \ldots$ cut out on $C$ a complete series (for $i$ high enough) now translates into the property of the system of the $\Phi^{\nu}$ of cutting out on a general plane \emph{all} the curves of order $\nu$ adjoint to the section of $F$, when $\nu \geqq n-3+i$; while for the values of $\nu < n-3+i$ the $\Phi^{\nu}$ cut out on the plane a part only of the adjoint curves, if $\delta_{i-1}, \delta_{i-2}, \ldots$ are greater than zero. Now the curves of order $\nu$ adjoint to the general plane section of $F$ form a system which has the dimension
\[
\varrho_{\nu} = \binom{\nu+2}{2} - 1 - k,
\]
where $k$ is a constant which depends only on the singularities of the section named, that is to say on the multiple curves of the $F$; $\nu$ may have any value $\geqq n-3$. It follows (by a very simple calculation) that the dimension of the system of the $\Phi^{\nu}$ for $\nu \geqq n-3+i$ is
\[
r_{\nu} = \binom{\nu+3}{3} - 1 - k\nu + k', \tag{1}
\]
where $k'$ is a new constant which depends on the multiple curves and on the isolated multiple points of $F$. The second member of the (1), for the values of $\nu < n-3+i$, no longer furnishes (in general) the exact, \emph{effective} dimension of the system of the $\Phi^{\nu}$ (for then one must attend to the defects $\delta_{i-1}, \delta_{i-2}, \ldots$); it furnishes only a character which one might call the \emph{virtual dimension} (numerical), and which coincides with the \emph{effective} dimension only in the case where $\delta_{i-1} = \delta_{i-2} = \cdots = 0$.

Precisely one finds, without difficulty, that the effective dimension of the system of the $\Phi^{n-4}$ is
\[
r_{n-4} = \binom{n-1}{3} - 1 - k(n-4) + k' + \sum_{0}^{i-1} \delta_{h},
\]

\textbf{*)} See the citations in n${}^{\text{o}}$ 4.

\origpage{286}
whence one deduces at once (having regard to the definition $\displaystyle p_{g} - p_{n} = \sum_{0}^{i-1} \delta_{h}$)
\[
p_{n} = \binom{n-1}{3} - k(n-4) + k'.
\]
This relation is stated by the theorem:

\emph{The numerical genus of a surface is nothing else than the virtual dimension of the system of the adjoint surfaces of order $n-4$, increased by one unit.}

It is precisely this dimension increased by 1 which, according to the cited authors, furnishes the definition of $p_{n}$.

On the subject of the preceding considerations it will not be useless to add some remarks.

We have arrived at the relation (1) without specifying the singularities (multiple curves and points) of the surface $F$ which one had to consider; $k$ and $k'$ depend on these singularities, it is true, but the expressions of $k$ and $k'$ as functions of them have no importance, at least in the theoretical considerations which we have just made. It would be otherwise if one wished actually to calculate $p_{n}$ for a given surface. But even in this case one might avoid the direct determination of $k$ and $k'$, if one succeeded in determining by some route the dimensions of the systems of the $\Phi^{\nu}$ corresponding to sufficiently high values of $\nu$; for we have just seen, in conclusion, that \emph{when $\nu$ runs through the increasing series of the whole numbers} (starting from a certain limit), \emph{these dimensions run through an arithmetic progression of the third order which, prolonged also to the lower values of $\nu$, contains $p_{n} - 1$ among its terms.}

At the same conclusion \emph{Cayley} and \emph{Nöther} also arrive directly, by calculating the expressions of $k$ and $k'$ by means of the singularities of the surface, and substituting them in the formula (1) (\emph{postulation formula})${}^{*)}$.

As soon as this calculation is made, a route presents itself for proving the invariance of $p_{n}$; it suffices in fact to make an analogous calculation over again with the characters of a surface transformed from $F$, and to compare the results. It is the route which the cited authors have followed. We have remarked above the difficulties which one meets in it; we have now just seen how one may avoid them.

\textbf{*)} See on this subject: \emph{Nöther}, \emph{Sulle curve multiple di superficie algebriche}, Annali di Matem.\ serie II, t.\ 5${}^{\text{o}}$, (1871) pag.\ 172.

\origpage{287}

\subsection*{27. Two new definitions of the numerical genus; the characteristic series of a linear system.}

The definition which we have given above of the character $p_{g} - p_{n}$, by the aid of a system $|C|$ and of the adjoint system $|C'|$, may even be presented in the following form.

The series which $|C'|$ cuts out on the general curve of $|C|$ may have, we have said, a certain defect $\delta_{0}$. This defect does not remain invariable in general when one replaces the system $|C|$ and its adjoint $|C'|$ by a new system $|D|$ and by its adjoint $|D'|$. Let us consider all the values which $\delta_{0}$ may take in correspondence with the different systems situated on the surface. It is proved that these values have a \emph{maximum}, which is naturally an invariant of the surface; but it is not a new invariant; on the contrary this maximum is precisely $p_{g} - p_{n}$. One has therefore the theorem${}^{*)}$:

\emph{The defect of the series cut out on the general curve of a linear system by the adjoint system may depend on the system which one chooses on the surface; but in all cases this defect attains a maximum, which is equal to the difference between the geometric and numerical genera of the surface.}

If one compares the definition of $p_{g} - p_{n}$ which thus results with that which we have given above, one sees that one may always choose on a surface a system $|D|$ such that the series cut out on the curve $D$ by $|D'+D|, |D'+2D| \ldots$ result complete${}^{**)}$- One may even prove that one arrives at such a system $|D|$ by taking a sufficiently high multiple of any system $|C|$ of the surface.

In the last definition of $p_{g} - p_{n}$ one attends to the series cut out by $|C'|$ on $C$. Now on the curve $C$ one finds another remarkable series, which we defined above (n${}^{\text{o}}$ 7); it is the \emph{characteristic series} $g_{n}^{r-1}$ cut out on the general curve of the system $|C|$ $\infty^{r}$ by the other curves of the system. We shall now occupy ourselves with this series; will it always be complete? or else, in the contrary case, what can be said of its defect?

It is naturally supposed here that the system $|C|$ is complete; for on the contrary hypothesis, if $|C|$ were contained in a system $\infty^{r+\delta}$ of the same degree $n$, the characteristic series $g_{n}^{r-1}$ of $|C|$ would be contained

\textbf{*)} \emph{Enriques}, \emph{Introduzione}... n${}^{\text{o}}$ 40. One may verify this theorem on the \emph{hyperelliptic} surfaces $(p_{g} = 1, p_{n} = -1)$; see on this subject \emph{Humbert}, \emph{Théorie générale des surfaces hyperelliptiques} n${}^{\text{o}}$ 130, 144; l.\ c.

\textbf{**)} \emph{Enriques}. l.\ c.; for the hyperelliptic surfaces see \emph{Humbert}, l.\ c., n${}^{\text{o}}$ 144. For the ruled surfaces see above, n${}^{\text{o}}$ 26.

\origpage{288}
in its turn in the characteristic series $g_{n}^{r+\delta-1}$ of the new system, and would consequently have a defect $\geqq \delta$. Now if $|C|$ is complete, the characteristic series may be complete, or else may not be so; this depends on the surface on which the system is situated. It is the first case that presents itself, for example, if the surface is rational; we remarked it above for the systems of plane curves (n${}^{\text{o}}$ 11), and it is naturally a question here of a property invariable with respect to the birational transformations. The same property is verified on the general surfaces of the order $n$, etc. Let us consider on the contrary an irrational ruled surface which is not a cone, for example the surface of the fourth order having two double directrix straight lines; the plane sections of this surface do indeed form a complete linear system, and yet the characteristic series $g_{4}^{2}$ is not complete; on the contrary it has the defect 1. What then is the bond which passes between the defect of the characteristic series and the invariants of a given surface?

If one considers several complete linear systems on one and the same surface, and attends to the defects of their characteristic series, one will sometimes find that these defects differ from one another. Let one take the maximum among the defects which correspond to all the systems of the surface; this maximum, which is finite and always exists (as may be proved), is naturally an invariant of the surface. Will it be a new invariant? One might think so. Well then, it is not so; for it is shown that this maximum is precisely equal to the difference $p_{g} - p_{n}$. One has therefore the following theorem${}^{*)}$:

\emph{The defect of the characteristic series of a complete system situated on a surface may depend on the system which one chooses on the surface; but in all cases this defect attains a maximum, which is equal to the difference between the geometric and numerical genera of the surface.}

\subsection*{28. The regular surfaces.}

The two preceding theorems give at once the fundamental properties of the surfaces which have $p_{g} = p_{n}$, surfaces which we call \emph{regular}. One has in fact:

a) \emph{On a regular surface the system adjoint to any (irreducible) linear system cuts out on the general curve of the latter the complete}

\textbf{*)} \emph{Castelnuovo}, \emph{Alcuni risultati}... n${}^{\text{o}}$ 7; there will be found the proof only for the case $p_{g} = p_{n}$; the Author proposes to set out the complete proof shortly. One may verify the theorem on the hyperelliptic surfaces (\emph{Humbert}, l.\ c., n${}^{\text{o}}$ 114), and on the ruled surfaces (\emph{Segre}, Mathem.\ Annalen, 34, n${}^{\text{o}}$ 3.).

\origpage{289}
\emph{canonical series;} it has the dimension $p_{g} + \pi - 1$, if $\pi$ is the genus of the last curve.

b) \emph{On a regular surface the characteristic series of every complete linear system is itself complete.}

Each of the two theorems may serve to characterize a regular surface. As regards the first, one may even say more, for \emph{in general the knowledge of a single system $\infty^{r} (r \geqq 3)$, on the curves of which the adjoint system cuts out the complete canonical series, is sufficient to affirm that the surface is regular}${}^{*)}$. The words \emph{in general} relate to the following restrictions which one meets in the proof: 1) the system has no proper fundamental curves; 2) the curves of the system which pass through a point of the surface do not in consequence pass through other points variable with that one. Are these necessary restrictions?

\subsection*{29. On some classes of irregular surfaces.}

The necessity of considering the two characters $p_{g}$ and $p_{n}$ of a surface depends on this, that there are numerous examples of surfaces (\emph{irregular}) having $p_{g} > p_{n}$. We have already cited in this connection the ruled surface with sections of genus $\pi$, for which one has $p_{g} = 0$, $p_{n} = -\pi$. But two very extensive families of irregular surfaces are now known.

The first family, which comprises also the ruled surfaces, is composed of the surfaces which contain an irrational pencil of curves, that is to say a non-linear system $\infty^{1}$ such that through every general point of the surface there passes a single curve of the system. \emph{The knowledge of an irrational pencil on a surface always suffices to affirm that the surface is irregular.} In other terms: \emph{on a regular surface every pencil of curves is a linear system}${}^{**)}$.

To construct the most general surface containing an irrational pencil, it suffices to transform an irrational ruled surface by the aid of a transformation $(1, n)$ rational in one sense only (expressing the coordinates of a point of the ruled surface by rational functions of the coordinates of a point of the transformed surface${}^{***)}$).

\textbf{*)} \emph{Castelnuovo}, \emph{Alcuni risultati}... n${}^{\text{o}}$ 6.

\textbf{**)} \emph{Castelnuovo}, \emph{Alcuni risultati}... n${}^{\text{o}}$ 10. The knowledge of an irrational pencil suffices also to recognize that the surface possesses integrals of total differentials of the first kind.

\textbf{***)} Examples of such surfaces will be found in a Note of M.\ \emph{Castelnuovo}, \emph{Osservazioni intorno alla geometria sopra una superficie}, Rendic.\ Istituto Lombardo, serie II, vol.\ 24, pag.\ 135.

\origpage{290}

A second family of irregular surfaces is obtained by considering the surfaces whose points correspond, element by element, to the couples of points of any curve of genus $p > 0$. Thus, if one starts from an elliptic curve $(p = 1)$, one arrives at the elliptic ruled surface $(p_{g} = 0, p_{n} = -1)$; if one starts from the curve of genus $p = 2$, one obtains a hyperelliptic surface $(p_{g} = 1, p_{n} = -1)$${}^{*)}$; etc.

There certainly exist other families of irregular surfaces, but up to the present none, so we believe, has formed the subject of a particular study.

\subsection*{30. On new invariant systems which belong to a surface.}

We have spoken up to the present of a single invariant system situated on a given surface, namely of the canonical system $|K|$. But one may at once introduce new invariant systems, for it suffices to consider the multiples $|K_{i}| = |iK|$ of the canonical system. It would not even be worth the trouble to attend to these new systems, if it did not sometimes happen that analogous systems are met with on surfaces which possess no canonical system $(p_{g} = 0)$. In this case, however, one cannot define $|K_{i}|$ as a multiple of $|K|$, for the latter no longer exists; but it suffices to have recourse to the definition of $|K|$, which extends at once. We have proved that if $|C|$ is any system on a surface, and $|C'|$ is its adjoint, the residual system $|K| = |C'-C|$ does not depend on the system $|C|$ from which one starts; likewise one sees directly that the system $|K_{i}| = |iC'-iC|$ does not depend on $|C|$ (up to exceptional curves), and may therefore also be obtained by subtracting the system $|iD|$ from the system $|iD'|$, where $|D'|$ is the adjoint of $|D|$${}^{**)}$. The system $|K_{i}|$ therefore enjoys the property of invariance; it will be called the \emph{i-canonical system}; $i$ is any positive whole number. Besides by the relation $|K_{i}| = |iC'-iC|$, this system is also defined by the relation $|K_{i}| = |C^{(i)}-C|$ (up to exceptional curves), where $C^{(i)}$ is the system which one obtains from $|C|$ by applying in succession $i$ times the operation of adjunction. The \emph{i-canonical} system furnishes us new invariant characters of the surface; such are for example the number $P_{i}$ of the i-canonical curves which are linearly independent, the genus of these curves, etc. For $i = 1$, $P_{i} = P_{1}$ gives us the geometric genus $p_{g}$ of the surface; for $i = 2$ one has a character $P_{2}$ which one might call the \emph{bigenus}, and which is to play (as will be seen) a

\textbf{*)} \emph{Humbert}, \emph{Théorie générale des surfaces hyperelliptiques}, Journal de Mathém., 4${}^{\text{e}}$ série, t${}^{\text{e}}$ IX; see also for $p = 3$ two Notes of the same author in the Comptes Rendus de l'Ac.\ d.\ Sc., Février 1895.

\textbf{**)} \emph{Enriques}, \emph{Introduzione}..., n${}^{\text{o}}$ 39.

\origpage{291}
fundamental role in the theory of the rational surfaces; etc. Some of these new characters may be expressed as functions of the old ones; not all, however; this is what we shall recognize on the subject of $P_{2}$.

But let us remark first that if a surface $F$ of the order $n$ is given (by its projective characters) in ordinary space, the i-canonical system on $F$ may be cut out by surfaces of the order $i(n-4)$ which behave in a \emph{suitable} manner along the multiple curves and at the multiple points of $F$. To specify the behaviour of the \emph{i-adjoint} surfaces in the general case would require many details relative to the higher singularities; let us be permitted to confine ourselves to the case $i = 2$, on the hypothesis that the surface $F$ possesses no other singularities than a double curve, and points of order $k = 3, 4, \ldots$ which are of order $\dfrac{k(k-1)}{2}$ for the double curve. \emph{In this case the bicanonical system is cut out on $F$ by the surfaces of order $2(n-4)$ which pass twice through the double curve of $F$.}

Let us return now to the relations between the canonical system $|K|$ and the i-canonical system $|K_{i}|$. If the first system exists, the second will likewise exist, whatever $i$ may be, for one has $|K_{i}| = |iK|$, and in consequence $P_{i} \geqq p_{g}$; therefore if $p_{g} > 0$, one has also $P_{i} > 0$. Suppose, on the contrary, that one has $p_{g} = 0$; $|K|$ no longer exists. Then there are two cases to be distinguished. It may in fact happen that $|K_{i}|$ does not exist either, whatever $i$ may be; this is what one may verify at once on the plane, on any ruled surface, more generally on every surface containing a system $|C|$ of genus $\pi$ two curves of which have more than $2\pi - 2$ variable intersections; on this hypothesis, in fact, the system $|iC|$ cannot be contained in the system $|iC'|$. But it may on the contrary happen that $|K_{i}|$ exists, in relation with some values of $i$, even when $|K|$ does not exist; in other terms, it may happen that one has $p_{g} = 0$, $P_{i} > 0 (i > 1)$. This is a singular property, which has no analogue in the theory of algebraic curves. That it sometimes presents itself is what we shall recognize on some examples, confining ourselves to the value $i = 2$.

A first example, very simple, is furnished by \emph{the surface of the sixth order which passes twice through the six edges of a tetrahedron}, and in consequence three times through the vertices of the latter${}^{*)}$; the surface has the geometric genus $p_{g} = 0$, for there exists no surface of the second order passing simply through the six edges. There exists on the contrary a bi-adjoint surface of the order 4, which is formed by the four faces of the tetrahedron; consequently one has $P_{2} = 1$. The last surface does not cut the given surface

\textbf{*)} \emph{Enriques}, \emph{Introduzione}..., n${}^{\text{o}}$ 39; \emph{Castelnuovo}, \emph{Sulle superficie di genere zero}, n${}^{\text{o}}$ 15.

\origpage{292}
outside the multiple lines; one will therefore say that the corresponding bicanonical curve has the order zero.

Here is a second example${}^{*)}$. Let us consider in space a straight line $d$, a conic section $k$ not meeting the straight line, and four points $A_{1}$, $A_{2}$, $A_{3}$, $A_{4}$ in general positions. Let us now construct a surface of the order 7 which passes three times through $d$ and twice through $k$, with the further condition that two sheets of the surface touch at $A_{i}$ the plane $A_{i}d$ $(i = 1, 2, 3, 4)$. It is proved first of all that such a surface exists; next that it has the genus $p_{g} = 0$ and the bigenus $P_{2} = 2$; finally that the bicanonical curves are $\infty^{1}$ curves of the fourth order and of genus 1, cut out by the planes passing through the straight line $d$.

\subsection*{31. On some relations between the invariant characters.}

The principal characters of a surface which we have just defined $(p_{g}, p_{n}, P_{2}, p^{(1)})$, although independent two by two, satisfy some inequalities which it is well to know. We have already remarked the two following, which flow from the definition,
\[ p_{g} \geqq p_{n}, \quad P_{2} \geqq p_{g}. \]
One may even say more, if one attends to the fact that the bicanonical system is the adjoint of the canonical system; in fact, if one supposes that the latter is irreducible, one obtains, when $p_{g} > 1$,
\[ p_{g} + p^{(1)} \geqq P_{2} \geqq p_{n} + p^{(1)}. \]

This relation, which becomes an equality for the regular surfaces, does not always subsist if $p_{g} = 1$, even if one has recourse to the virtual definition of $p^{(1)}$. Precisely, confining ourselves to the case $p_{g} = p_{n} = 1$${}^{**)}$, one finds
\[ P_{2} = p^{(1)} + 1, \quad \text{or} \quad P_{2} = p^{(1)}, \]
according as there exists, or does not exist, a canonical curve of order $> 0$ (outside the exceptional curves).

Consequently the minimum of $p^{(1)}$ for the regular surfaces of genus 1 is 1. There are two kinds of such surfaces for which the minimum is attained; they correspond to the values $P_{2} = 2$, $P_{2} = 1$${}^{***)}$. The surfaces of the second kind are especially remarkable. They may be represented on a surface $F$, of a certain order $n$, possessing

\textbf{*)} \emph{Castelnuovo}, l.\ c.

\textbf{**)} \emph{See} for this case \emph{Enriques}, \emph{Sui piani doppi di genere uno}, n${}^{\text{os}}$ 4, 5.

\textbf{***)} An example of surfaces of the first kind is furnished by the surface of the fifth order of n${}^{\text{o}}$ 23. As for the surfaces of the second kind, one has first of all the general surface of the fourth order, next the surface of the sixth order which passes doubly through the intersection of two surfaces of the order 2 and 3; etc.

\origpage{293}
no exceptional curve; the $F$ is not met by the adjoint surface of the order $n-4$ outside the multiple curves. On $F$ every complete linear system $|C|$ which has no base points enjoys a remarkable property; \emph{if $\pi$ is the genus of $|C|$, the dimension of $|C|$ has the same value $\pi$, and two curves $C$ meet in $2\pi-2$ points $(r=\pi, n=2\pi-2)$. The system $|C|$ is the adjoint of itself.}

Let us return now to the relations of inequality between the characters of a surface, and consider $p_{g}$ and $p^{(1)}$, supposing, for simplicity, that the canonical system is irreducible $(p_{g} > 1)$.

M.\ \emph{Nöther}${}^{*)}$ has proved that one always has
\[ p^{(1)} \geqq 2p_{g} - 3. \]
\emph{If the minimum of $p^{(1)}$ is attained, the canonical curves are hyperelliptic} (and \emph{viceversa}); they meet two by two in
\[ \frac{p^{(1)} - 1}{2} = p_{g} - 2 \]
couples of conjugate points. Consequently if $p_{g} > 2$, one may take as the equation of the surface (or of a transform of it)
\[ z^{2} = f(x, y), \]
where $f$ is a polynomial. One may even be more precise${}^{**)}$, for \emph{one may arrange matters in such a way that}

1) \emph{either $f$ has the degree 6 in $x$ (and the total degree $2n > 6$)},

2) \emph{or else $f$ has the total degree 8 or 10 (respect.\ $p_{g} = 3, 6$)}.

In case 1) the planes $y = const.$ cut out on the surface a pencil of curves of the genus 2.

Let one now set aside the surfaces for which
\[ p^{(1)} = 2p_{g} - 3; \]
it is then proved that $p^{(1)}$ cannot descend below $3p - 6$, namely
\[ p^{(1)} \geqq 3p_{g} - 6; \]
and precisely${}^{***)}$:

\emph{Every surface having}
\[ p^{(1)} = 3p_{g} - 6 \qquad (p_{g} > 3) \]
\emph{may be represented by an equation}
\[ f(x, y, z) = 0 \]
\emph{of the degree 4 with respect to $x$, $y$ (or by a transform of it); save in the case of certain surfaces of the genus $p_{g} = 4$ (the general surface}

\textbf{*)} \emph{Zur Theorie}... Math.\ Ann.\ 8, p.\ 252.

\textbf{**)} \emph{Enriques}, \emph{Sopra le superficie algebriche di cui le curve canoniche sono iperellittiche}. Rendic.\ Accad.\ Lincei 1896.

\textbf{***)} \emph{Castelnuovo}, \emph{Osservazioni interno alla geometria sopra una superficie}, Nota II, Rendic.\ Istit.\ Lombardo 1896.

\origpage{294}
\emph{of the fifth order), $p_{g} = 5$, $p_{g} = 7$.} In the general case the planes $z = const.$ cut out on the surface a pencil of curves of the genus 3 (and of order 4).

\section*{Chapter VI. Special and non-special linear systems. Extension to surfaces of the theorem of Riemann-Roch.}

Before leaving the general theory of surfaces we must devote a few words to the questions which relate to the extension of the theorem of Riemann-Roch to surfaces.

But first let us examine what is the true nature of this theorem when one confines oneself to curves.

\subsection*{32. Special and non-special series on a curve.}

With respect to the canonical series $g_{2p-2}^{p-1}$ belonging to a curve of genus $p$, the series situated on the same curve are distributed into two families${}^{*)}$. To the first family belong the series which are contained in the canonical series, and the latter itself; they are called \emph{special series}; the special series have naturally the order $\leqq 2p - 2$ and the dimension $\leqq p - 1$. The other family is formed by the series which are not contained in the canonical series; these are the \emph{non-special series}, whose order and dimension have no upper limit.

If one wishes to see clearly the difference which passes between these two families, it is fitting to have recourse to the complete series.

\emph{Let $g_{n}^{r}$ be a complete series on a curve of the genus $p$; then}

1) \emph{if $g_{n}^{r}$ is non-special one has the relation}
\[ r = n - p, \]
which permits one of the three characters $n$, $r$, $p$ to be calculated by the aid of the others;

2) \emph{if, on the contrary, $g_{n}^{r}$ is special, the preceding equality is to be replaced by the inequality}
\[ r > n - p. \]
And, inversely, if the last inequality subsists between the characters of any series (complete or incomplete), one may affirm that the series is special.

The distinction which we have just made may even be presented in the following form. Let us regard $n - p$ as the \emph{virtual dimension} (theoretical) of a complete series of order $n$ on a curve of genus $p$,

\textbf{*)} \emph{Brill} and \emph{Nöther}, \emph{Ueber die algebraischen Functionen}..., Math.\ Ann.\ 7.

\origpage{295}
and let us compare the virtual dimension with the \emph{effective dimension} $r$ which the series actually possesses. Then we see that if the series is non-special the two dimensions coincide, and the series is, so to say, \emph{regular}; if, on the contrary, the series is special the effective dimension is higher than the virtual dimension, and the series is irregular.

For a complete special series $g_{n}^{r}$ there is occasion to consider the difference between the two dimensions
\[ i = r - (n-p) = p - n + r, \]
It is proved that \emph{$i$ is precisely the number of linearly independent canonical groups which contain a general group of the series $g_{n}^{r}$}. This is the well known theorem which bears the name of \emph{Riemann-Roch}. The number $i$ is a character of the series $g_{n}^{r}$ which one might call the \emph{degree of speciality}; for a non-special series one has $i = 0$, for the canonical series $g_{2p-2}^{p-1}$, $i = 1$; etc.

Profiting by the notion of the residual series (n${}^{\text{o}}$ 9), one may state the theorem of \emph{Riemann-Roch} in the following form, which is due to MM.\ \emph{Brill} and \emph{Nöther}, and which lends itself best to geometric applications:

\emph{Every complete special series $g_{n}^{r}$ situated on a curve of genus $p$ has, with respect to the canonical series, a residual series $g_{n'}^{r'}$ whose characters have the following values}
\[ n' = 2p - 2 - n, \quad r' = p - 1 - (n-r) = i - 1. \]

\subsection*{33. Some remarks on the linear systems of plane curves.}

The distribution of the series into two families extends at once to the linear systems of curves on a surface, and gives rise to the notion of \emph{special systems} and \emph{non-special systems}, of which the first are contained in the canonical system (and have consequently the dimension $\leqq p_{g} - 1$, the genus $\leqq p^{(1)}, \ldots$), while the second are not contained in it. Now one might expect to find for the non-special systems a \emph{regularity} analogous to that which we met with in the non-special series. But it is quite otherwise. This is what we shall recognize first of all in the simplest case, on the \emph{plane} $(p_{g} = 0)$, where all the systems are non-special.

A complete system $|C|$ of plane curves of order $m$ is entirely defined as soon as one has fixed the base points of the system, and the orders of multiplicity which the curves $C$ have at these points. Now to calculate the dimension of $|C|$, at least with some approximation, one may first evaluate the number of the conditions which each base point, separated from the others, presents to the curves of order $m$; one will add

\origpage{296}
the numbers which correspond to all the base points, and will subtract the sum from the number $\frac{1}{2} m(m+3)$, which expresses the dimension of the system formed by all the curves of the order $m$. Let us designate by $\varrho$ the result of this calculation; will $\varrho$ furnish us exactly the dimension $r$ of $|C|$? Certainly if the base points present conditions all independent to the curves of order $m$; this is what happens, for example, if there are no base points, or if these are found at general points of the plane; one has then the relation $r = \varrho$. But it is no longer so if $\omega$ among the conditions imposed by the base points are a consequence of the others; for one finds then $r = \varrho + \omega$. An example of this case is furnished us by the system of the curves of order $m$ which pass through the $m\mu$ intersections of one of them with a general curve of order $\mu < m$; one has in fact
\[ r = \frac{(m-\mu)(m-\mu+3)}{2} + 1, \]
\[ \varrho = \frac{m(m+3)}{2} - m\mu, \]
and, consequently,
\[ \omega = r - \varrho = \frac{(\mu-1)(\mu-2)}{2}. \]

Although the number $\varrho$ may, in many cases, differ from the exact dimension $r$ of a linear system, it is nevertheless fitting to attend to it, for it is a new invariant of the system (with respect to the birational transformations). We call $\varrho$ the \emph{virtual dimension} of the system $|C|$, $r$ the \emph{effective dimension}. In consequence we distribute the complete linear systems of plane curves (which are all non-special) into two families${}^{*)}$:

1) the first family comprises the \emph{regular systems}, for which the effective dimension is equal to the virtual dimension, $r = \varrho$; they correspond, in a certain fashion, to the non-special (or regular) series on a curve;

2) the second family comprises the \emph{surabundant} systems, whose effective dimension exceeds the virtual dimension, $r = \varrho + \omega$; $\omega > 0$ is the \emph{surabundance} of the system. It is not fitting to regard these systems as the analogues of the special series; on the contrary, they differ from them in many properties. For example, in this, that on a curve of genus zero there are no special series, while on the plane $(p_{g} = p_{n} = 0)$ we find surabundant systems; in this, again, that the order and the dimension of a special series cannot exceed certain limits, while the degree, the dimension, the genus

\textbf{*)} \emph{Castelnuovo}, \emph{Ricerche generali sopra i sistemi lineari di curve piane}, l.\ c.\ n${}^{\text{o}}$ 1.

\origpage{297}
of a surabundant system may be as high as one desires. It is fitting rather to regard the \emph{surabundance} as a \emph{cause of irregularity} of another nature than the \emph{speciality} considered on curves, although it has the same effect of increasing the effective dimension with respect to the virtual dimension.

The distinction which we have made between the regular and the surabundant systems translates into a distinction between the characteristic series of these systems. One has in fact the theorem${}^{*)}$:

\emph{A system is regular or else surabundant according as the characteristic series of the system is non-special or else special;} whence it follows, for example, that a surabundant system of curves of genus $\pi$ has the degree $\leqq 2\pi - 2$ and the effective dimension $\leqq \pi - 1$${}^{**)}$.

In general if one designates (as ordinarily) by $n$ and $r$ the degree and the effective dimension of a complete system, one has for a regular system the equality
\[ \pi - n + r - 1 = 0, \tag{1} \]
while for the surabundant systems one has
\[ \pi - n + r - 1 = \omega > 0 \tag{2} \]
\[ (\omega \text{ being the surabundance}); \]
in all cases
\[ \varrho = n - \pi + 1 \tag{3} \]
is the virtual dimension of the system.

\subsection*{34. Regular systems on a surface.}

The distinction which we have just made between the \emph{regular} and \emph{surabundant} systems on the plane (or on the rational surfaces) may now be carried over to any surface. Here one meets, however, a complication which depends on the distribution of the systems into special and non-special, of which we have spoken above. To avoid this difficulty we confine ourselves first of all to the \emph{non-special} systems (namely, to the systems which are not contained in the canonical system).

Let $|C|$ be a non-special, complete system, whose characters we shall designate by $n$, $r$, $\pi$, according to custom. Let us consider the expression
\[ \pi - n + r - 1 \]
which we have already considered on the plane. It may take different

\textbf{*)} \emph{Castelnuovo}, l.\ c., n${}^{\text{o}}$ 18.

\textbf{**)} \emph{Segre}, \emph{Sui sistemi lineari}, Circolo Matem.\ di Palermo, t.\ I, 1887.

\origpage{298}
values, according to the system $|C|$ to which it relates; but it is proved${}^{*)}$ that among all these values there is \emph{a minimum, which is precisely the numerical genus $p_{n}$ of the surface} (for the plane one has $p_{n} = 0$). Well then, we call \emph{regular} every (non-special) system whose characters satisfy the relation
\[ \pi - n + r - 1 = p_{n}; \tag{4} \]
\emph{surabundant} every system for which one has
\[ \pi - n + r - 1 = p_{n} + \omega, \tag{5} \]
where $\omega > 0$ designates the \emph{surabundance} of the system. We shall consider first of all the regular systems, for they are the simplest, and they satisfy very remarkable general laws. Here are some of them${}^{**)}$.

\emph{The characteristic series of a complete regular system has the defect $p_{g} - p_{n}$}, that is to say the maximum defect which belongs to complete systems (n${}^{\text{o}}$ 27).

\emph{The canonical system cuts out on the general curve of a regular system a complete series}, whose dimension is $p_{g} - 1$, and whose order is $2\pi - 2 - n$; it is the series residual to the characteristic series with respect to the canonical series.

In particular: \emph{if the surface has the geometric genus $p_{g} = 0$, the characteristic series of a regular system is non-special}; this is what happens for example on the plane, or on the ruled surfaces.

One may construct a regular system whose dimension is as great as one wishes; to this end one will begin by constructing in ordinary space a surface $F$ in birational correspondence with the given surface, and having no \emph{isolated} multiple points (n${}^{\text{o}}$ 15); one will then draw the surfaces adjoint to $F$. It is proved that these surfaces cut out on $F$ a regular system as soon as their order is high enough. If, for example, the $F$ were a general surface of the order $n$, one would easily see that the general surfaces of any order $> n - 4$ cut out on $F$ a regular non-special system.

The definition of the \emph{surabundance} $\omega$ which we have given is altogether arithmetical; we are now in a position to recognize its geometric signification and to justify its name.

But we must first introduce a new character of any linear system $|C|$ (special or non-special) on the surface;

\textbf{*)} For the regular surfaces this theorem will be found proved in the \emph{Ricerche di geometria}... of M.\ \emph{Enriques}, IV, 2; for any surface see \emph{Castelnuovo}, \emph{Alcuni risultati}..., n${}^{\text{o}}$ 8. It is to be remarked that this theorem furnishes a new definition of $p_{n}$.

\textbf{**)} \emph{Castelnuovo}, l.\ c.

\origpage{299}
a character which becomes the \emph{virtual dimension} of the system (n${}^{\text{o}}$ 33) as soon as the surface is rational, and which will consequently be designated in all cases by the same name.

Let us construct to this end a \emph{regular} system $|D|$ wide enough to contain $|C|$; this is what we may do in several ways. To deduce the system $|C|$ from the new system $|D|$ it will in general be necessary to subtract from $|D|$ a suitable curve $\gamma$, and to impose certain base points on the residual system $|D-\gamma|$. There now presents itself a means for calculating the dimension $r$ of $|C|$; for it suffices to subtract from the dimension of $|D|$ the number of the conditions which must be imposed on the curves $D$ in order that they should contain the curve $\gamma$ and the base points named. Now if we confine ourselves first to a first approximation, we may suppose: 1) that the series cut out by $|D|$ on the curve $\gamma$ is complete and non-special; 2) that the conditions imposed by the base points named are all independent among themselves and of the preceding ones. In executing the calculation on these hypotheses we arrive at a result $\varrho$ which does not generally coincide with the true, \emph{effective} dimension $r$ of $|C|$, but which nevertheless does not lack interest. It is proved in fact that one arrives at the same result $\varrho$ if, on the same hypotheses, one calculates the dimension of $|C|$ starting from any regular system other than $|D|$ which contains $|C|$. Consequently $\varrho$ is a true character of the system $|C|$; we shall call it \emph{the virtual dimension} of $|C|$. It is expressed, moreover, easily by the aid of the other characters $n$, $\pi$ of $|C|$, and of $p_{n}$, for one has
\[ \varrho = p_{n} + n - \pi + 1, \tag{6} \]
a relation which comprises, in particular, the (3) of n${}^{\text{o}}$ 33.

This laid down, let us now return to our non-special system $|C|$, whose effective and virtual dimensions we designate by $r$ and $\varrho$. If we compare these characters by the aid of the relations (4), (5) and (6), we find at once that if $|C|$ is regular one has $\varrho = r$; if, on the contrary, $|C|$ is surabundant, with the surabundance $\omega$, then $\varrho < r$, and precisely
\[ r - \varrho = \omega. \]

Therefore \emph{the surabundance of a non-special system is the difference between the effective dimension and the virtual dimension of the system}, in agreement with what we have recognized on the plane.

Another interpretation, this one altogether geometric, may be given of the character $\omega$ on the regular surfaces $(p_{g} = p_{n} = p)$. One arrives at it by considering the group of the $n$ variable intersections of two general curves of the given system $|C|$, and the curves of the adjoint system $|C'|$ which pass through this group. In fact, while the number of these last curves which are linearly independent is $2p$

\origpage{300}
if the system $|C|$ is regular, the same number is $2p + \omega$ if the system $|C|$ has the surabundance $\omega$. Consequently${}^{*)}$:

\emph{The surabundance of a non-special system $|C|$ on a regular surface of genus (geom.\ $=$ num.) $p$ is the excess over $2p$ of the number of the linearly independent adjoint curves $C'$ which pass through the variable intersections of two $C$.} One will verify the theorem at once on the plane.

\emph{Remark.} --- Suppose that the surface $F$ (regular or not) is given in ordinary space, and possesses only ordinary singularities. According to the \emph{Restsatz} (n${}^{\text{o}}$ 12), every system $|C|$ on $F$ may be constructed by the aid of the adjoint (or sub-adjoint) surfaces, of a sufficiently high order, which pass through a fixed curve $\gamma$ and through the base points of $|C|$. Then the \emph{virtual dimension} $\varrho$ of $|C|$ (and consequently its surabundance $\omega$) may be calculated according to the \emph{postulation formulas} of M.\ \emph{Nöther} (n${}^{\text{o}}$ 26).

\subsection*{35. Special systems.}

The properties which we have just stated on the non-special systems extend also to the special systems, provided one has regard to a new character which presents itself in this case, and which is altogether analogous to the character $i$ which we introduced in the study of the special series. A complete special system $|C|$ is contained, by definition, in the canonical system $|K|$; there is therefore occasion to consider the number $i$ of the linearly independent canonical curves which contain the general curve of $|C|$; this number $i$, which designates, so to say, the \emph{degree of speciality} of $|C|$, is the new character. If $|C|$ is non-special, one has $i = 0$; if $|C|$ is the canonical system, then $i = 1$; etc. In place of $i$ one may also consider the dimension $r_{1}$ of the system $|C_{1}| = |K-C|$ residual to $|C|$ with respect to the canonical system; one has naturally
\[ r_{1} = i - 1. \]

By the aid of the character $i$ we may extend in the following fashion the formulas which we have found apropos of the non-special systems. Let us designate by $n$, $\pi$, $r$, $i$ the characters of a special system $|C|$, and by $p_{n}$ the numerical genus of the surface; one has then
\[ \pi - n + r - 1 \geqq p_{n} - i, \]
and precisely if
\[ \pi - n + r - 1 = p_{n} - i + \omega, \tag{7} \]
one calls $\omega \geqq 0$ the \emph{surabundance} of the system; for $i = 0$ one falls back

\textbf{*)} \emph{Enriques}, \emph{Ricerche di geometria}... IV.\ 2.

\origpage{301}
into the formula of the non-special systems. One sees from this that \emph{a system is certainly special if}
\[ \pi - n + r - 1 < p_{n}; \]
but, because of $\omega$, one cannot say that this relation is also necessary; contrary to what happened on curves, where the inequality $n - r < p$ could serve as the definition of a complete special series.

One may write the equality (7) in such a way as to bring into evidence the dimension $r$ of the system $|K-C|$. One arrives at the following result:

\emph{On a surface of numerical genus $p_{n}$, every complete special system having the characters $n$, $r$, $\pi$ has, with respect to the canonical system, a residual system whose dimension is}
\[ r_{1} = p_{n} - (\pi-n+r) + \omega, \tag{8} \]
\emph{where $\omega$ (the surabundance of the first system) designates a number $\geqq 0$.}

This theorem must be regarded as the extension to surfaces of the theorem of \emph{Riemann-Roch} stated for curves, for the latter also gave us the dimension $r_{1}$ of the series residual to a series $g_{n}^{r}$ with respect to the canonical series $(r_{1} = p - 1 - (n-r))$. For surfaces we see that the formula is complicated by the presence of the number $\omega$, which has no analogue in the theory of curves${}^{*)}$${}^{**)}$.

The equality (7) may further be transformed by introducing the virtual dimension $\varrho$ of the system $|C|$, which we defined above for any system by the aid of the relation (6) of n${}^{\text{o}}$ 34. One finds thus the relation
\[ r - \varrho = \omega - i, \]

\textbf{*)} The extension of the theorem of R.-R.\ to surfaces is due to M.\ \emph{Nöther} (Comptes Rendus de l'Ac.\ des Sc.\ 1886), on the understood hypothesis that the characteristic series of the system $|C|$ is complete; we know at present that this hypothesis is always verified on a regular surface $(p_{g}=p_{n}=p)$. M.\ \emph{Nöther} presents the relation of the theorem under the form
\[ r_{1} \geqq p + n - \pi - r; \]
it is M.\ \emph{Enriques} (\emph{Ricerche di geometria}... IV, 2) who showed the geometric signification of the difference $\omega$ between the two members of the inequality; see on this subject the lines which follow above.

\textbf{**)} Besides the dimension $r_{1}$ of the system $|K-C|$, residual to $|C|$ with respect to the canonical system, one may also express the other characters $n_{1}$ (degree), $\pi_{1}$ (genus), $\omega_{1}$ (surabundance) of the new system by the aid of the characters of the system $|C|$. One has only to have recourse to the relations of n${}^{\text{o}}$ 9. One will thus find
\[ \begin{aligned} \pi_{1} &= p^{(1)} - 3(\pi-1) + 2n, \\ n_{1} &= p^{(1)} - 1 - 4(\pi+1) + 3n, \\ \omega_{1} &= \omega, \end{aligned} \]
where $p^{(1)}$ is the genus of the canonical system (\emph{Curvengeschlecht}). See on this subject \emph{Enriques}, l.\ c.

\origpage{302}
which tells us that \emph{for a special system the surabundance is obtained by adding the character $i$ to the difference between the effective and virtual dimensions of the system.} This definition of $\omega$ reduces to that which we set out on the subject of the non-special systems when $i = 0$.

For the regular surfaces $(p_{g} = p_{n} = p)$ one may even extend to the special systems the second definition of the surabundance which we stated for the non-special systems. One has only to replace, in the latter, $\omega$ by $\omega - i$.

\subsection*{36. Survey of the preceding considerations.}

Having regard to the complications which a linear system of curves on an algebraic surface may present, it will not seem useless that we should here briefly sum up the remarks which we have just made on this subject.

If $n$ and $\pi$ designate the degree and the genus of a complete system $|C|$ on a surface of numerical genus $p_{n}$, one must attend first to the virtual dimension of $|C|$, which is furnished by the expression
\[ \varrho = p_{n} + n - \pi + 1. \]

Next one must have regard to two causes of irregularity whose effect is to modify the effective dimension $r$ of $|C|$ with respect to the virtual dimension $\varrho$. These two causes have moreover opposite effects, which may sometimes cancel each other. They may act either separately or at the same time. They are:

1) the \emph{surabundance}, which makes the dimension of the system \emph{increase} $(r = \varrho + \omega)$; it may present itself on any surface, and for systems of dimension, of genus, and of order as high as one wishes; nothing analogous takes place for the linear series on a curve;

2) the \emph{speciality}, to which is due, on the contrary, a \emph{diminution} of the dimension $(r = \varrho - i)$; a property the more remarkable in that in the theory of linear series on a curve the analogous cause has the effect of increasing the dimension $(r = n - \pi + i)$. On surfaces, moreover (as on curves), there is occasion to consider this cause of irregularity only for the systems having their characters $n$, $\pi$ and $r$ lower than certain limits (which are respectively $p^{(2)}$, $p^{(1)}$ and $p_{g} - 1$).

The only systems on which neither the one nor the other of these causes acts are the \emph{regular systems} $(r = \varrho, \omega = i = 0)$.

\origpage{303}

\section*{Chapter VII. On the rational surfaces and on some double planes.}

We devote this last Chapter of our Monograph to the study of the class of the rational surfaces; a class remarkable also because the researches to which it has given rise have led, by natural extension, to researches on general algebraic surfaces. We shall then say a few words on the double planes.

The theory of rational surfaces began with the study of some surfaces, particular from the projective point of view, which could be represented birationally on the plane; one sought to read on the representation the principal properties of such surfaces. Thus one finds considered first of all the surface of the second order (\emph{Plücker}, 1847; \emph{Cayley} and \emph{Chasles}, 1861), and some years later the general surface of the third order (\emph{Cremona}, \emph{Clebsch}, 1865); there follow the ruled surface of the third order, the surface of \emph{Steiner}, etc.

Having recognized on these examples the profit that might be drawn from the representation of a surface for the study of the latter, having learnt to overcome the difficulties which the reading of the properties of the surface on the representation might present at first sight, the (projective) problem of studying a surface from its known representation might be regarded as resolved in advance in all cases. Nevertheless another problem (this one belonging to the \emph{Geometry on a surface}) presented itself: to determine whether a surface given by some one of its characters is rational, or else is not.

It is to \emph{Clebsch} and to M.\ \emph{Nöther} that the fundamental propositions on this subject are due. \emph{Clebsch}${}^{*)}$ showed how one might establish the rational character of a great number of surfaces by first representing these surfaces on \emph{a double plane} (see n${}^{\text{o}}$ 38), and then seeking the conditions of rationality of a double plane. M.\ \emph{Nöther}${}^{**)}$ completed these researches, and gave the capital theorem${}^{***)}$ according to which every surface containing a linear pencil of rational curves is rational.

One succeeded later in extending these propositions, as well as in adding others of a very general character, which we shall set out in the following pages. But it must nevertheless be remarked that most of the new results rest in conclusion on the theorems of \emph{Clebsch} and \emph{Nöther}.

\textbf{*)} \emph{Ueber den Zusammenhang}... Math.\ Annalen, 3, 1870.

\textbf{**)} \emph{Ueber die ein-zweideutigen Ebenentransformationen}, Sitzungsber.\ d.\ phys.\ medicin.\ Soc.\ zu Erlangen, 1878; \emph{Ueber eine Classe}..., Math.\ Annalen, 33.

\textbf{***)} \emph{Ueber Flächen, welche Schaaren rationaler Curven besitzen}. Math.\ Ann., 3, 1870.

\origpage{304}

\subsection*{37. Some properties of the rational surfaces.}

The general properties of surfaces which we have just studied in the preceding Chapters furnish properties of the rational surfaces as soon as one particularizes the genera of the surface $(p_{g} = p_{n} = 0)$. These same properties, moreover, may be deduced from the direct study of the systems of plane curves; and the last route is naturally preferable from the point of view of simplicity.

In one way or another one will arrive at the following propositions, which we bring together here, although some of them have already been stated.

\emph{On a rational surface:}

a) \emph{the characteristic series of every complete linear system is complete;}

b) \emph{the (canonical) series cut out on the general curve of a linear system by the adjoint system is itself complete.}

Each of these two theorems amounts to saying that \emph{a rational surface has the two genera $p_{g}$, $p_{n}$ equal}; each may be deduced from the other.

c) \emph{On a rational surface there exist linear systems whose characteristic series is non-special, and among these there are systems of genus $\pi$ whose degree $n > 2\pi - 2$;} such is for example the system formed by all the plane curves of a given order.

The theorem c) tells us that \emph{the geometric genus $p_{g}$ $(= p_{n})$ of a rational surface is null.} It tells us more: not only is there no canonical system on the surface, but neither the bicanonical system nor the i-canonical system ($i$ arbitrary) exists either; so that one has (adopting the same designations as in n${}^{\text{o}}$ 30)
\[ p_{g} = 0, \quad P_{2} = 0, \quad P_{3} = 0 \ldots . \]

We shall soon see that these conditions are not, moreover, all independent.

d) \emph{On a rational surface there exist linear systems of any genus $\pi$.} For $\pi = 0$ one finds systems having the dimension $r$, and the degree $n = r - 1$, as high as one wishes; it is no longer so for $\pi > 0$, for \emph{a system of curves of genus 1 has the dimension and the degree $\leqq 9$; while for a system of genus $\pi > 1$ the dimension cannot exceed $3\pi + 5$, nor can the degree exceed $4\pi + 4$;} but these \emph{maxima} may be attained${}^{*)}$.

\textbf{*)} For $\pi = 1$ the theorem depends on the reduction of the systems of plane elliptic curves to well known types; see the works of MM.\ \emph{Bertini} (Annali di Matematica, s${}^{\text{e}}$ 2${}^{\text{a}}$.\ t.\ 8, 1877), \emph{Guccia} (1886), \emph{Martinetti} (1887) (Rendic.\ Circolo Mat.\ di Palermo, t.\ I), \emph{Jung} (Rendic.\ Istituto Lombardo 1887---88). For $\pi > 1$ the upper limits of $n$ and $r$ were assigned by M.\ \emph{Castelnuovo} (\emph{Massima dimensione}..., Annali di Matem.\ s${}^{\text{e}}$ 2${}^{\text{a}}$, t.\ 18, 1890).

\origpage{305}

Without insisting further on the properties of the rational surfaces, we shall now approach a very important question which is, in a certain sort, the inverse of the preceding ones: it is a matter of deciding whether a given surface is rational, by the examination of some known characters of the surface. The problem, thus posed, is not well defined; it will be necessary to particularize the characters of the surface which one regards as known; this is what we shall do in several ways.

\subsection*{38. The degree of a linear system of curves situated on the surface is known.}

Often it is the knowledge of a linear system of curves on the surface that permits the rationality of the latter to be decided; sometimes even (in the simplest cases) the knowledge of a single character of the system (for example, of the degree $n$, or of the genus $\pi$) suffices; according to the character to which one attends, one arrives at different classes of results.

Let us suppose, first of all, that there exists on the surface a system $|C|$ whose \emph{degree} is $n > 0$; the dimension $r$ of the system will be $\geqq 2$ and $\leqq n + 1$. Let us examine successively the first values which may be attributed to $n$.

1) The first cases which we meet correspond to $n = 1$, $r = 2$, and $n = 2$, $r = 3$; these lead to rational surfaces, as one sees at once.

2) Let us stop rather on the more interesting case where one has $n = 2$, $r = 2$. If we make correspond to each curve $C$ of the system a straight line of a plane, we shall represent two conjugate points of the surface (intersections of two $C$) on a general point of the plane; the surface in this fashion is represented on a \emph{double plane}, that is to say on a surface whose equation may be written in the form $z^{2} = R(x, y)$, $R$ being a rational function (which one may suppose entire). It is therefore the question of the rationality of a double plane that we now meet.

A double plane is entirely defined by its \emph{branch curve} (\emph{Uebergangscurve}) $R(x, y) = 0$, which is the locus of a point whose two corresponding points on the surface coincide. What peculiarity must the branch curve present in order that the double plane should be rational (representable on the simple plane)? The question, approached by

\origpage{306}
\emph{Clebsch}, was entirely resolved by M.\ \emph{Nöther}${}^{*)}$. Here is the answer:

\emph{In order that a double plane should be rational, it is necessary and sufficient that the branch curve can be brought back, by a birational transformation of the plane, to one of the three curves which follow:}

a) \emph{a curve of a certain order $n$ having a multiple point of order $n - 2$;}

b) \emph{a curve of the fourth order;}

c) \emph{a curve of the sixth order having two infinitely near triple points.}

The degenerations of these three cases are admitted, except the splitting a${}_{1}$) of the curve of order $n$ into $n$ straight lines of a pencil, b${}_{1}$) of the curve of the fourth order into four straight lines of a pencil, c${}_{1}$) of the curve of the sixth order into six straight lines of a pencil, or into three conics of a pencil with two fixed contacts; for these degenerations lead to double planes which represent ruled surfaces${}^{**)}$.

Let us come now to the values $> 2$ of the degree $n$, but let us suppose at the same time $r > 2$, for up to the present nothing is known about the rationality of the triple planes, $\ldots$ $(n = 3, \ldots, r = 2)$. Let us suppose more precisely that one may construct a surface $F$ of order $n$ of ordinary space whose plane sections belong to the given system $|C|$, which has the degree $n$ and the dimension $r$.

It is a matter of examining in what cases the surface $F$ is rational. The question has been entirely resolved for the smallest values of $n(> 2)$.

3) One knows for example that \emph{every surface of the order $n = 3$ is rational, abstraction made of the general cone of the third order.}

4) One knows also that \emph{a surface of the order $n = 4$ having multiple lines} (at least one double straight line) \emph{is rational, except the ruled surfaces of genus $> 0$. Among the surfaces of the fourth order which have no multiple curves there are only four families of rational surfaces, which were determined by} M.\ \emph{Nöther}${}^{***)}$. One of them is composed of the surfaces of the fourth order having a triple point.

\textbf{*)} Memoirs cited.

\textbf{**)} The correspondence which passes between the points of a rational double plane and the couples of points of a simple plane gives rise, on the latter, to an \emph{involution of} $(\infty^{2})$ \emph{couples of points}. Now the types of these involutions which are distinct with respect to the birational transformations were determined by M.\ \emph{Bertini} (\emph{Ricerche sulle trasformazioni univoche involutorie del piano}. Annali di Matem.\ s${}^{\text{e}}$ 2${}^{\text{a}}$, t.\ 8, 1877); they correspond naturally to the three types of double planes. The possibility of the existence of other types besides the preceding, which was not excluded by the Memoir of M.\ \emph{Bertini}, was set aside by ulterior researches (\emph{Kantor}, \emph{Castelnuovo}).

\textbf{***)} \emph{Ueber die rationalen Flächen vierter Ordnung}, Math.\ Annalen, 33, 1888.

\origpage{307}
The three other families correspond to a singular double point which the surface possesses, from which the surface is projected on one of the three rational double planes which we have just cited; the first of these remarkable cases was pointed out by M.\ \emph{Cremona}.

5) For $n \geqq 5$ one has only incomplete results; that is why we do not insist on them for the moment, reserving to ourselves to indicate how the determination of the rational surfaces of a given order might now be approached.

\subsection*{39. The genus of a linear system of curves situated on the surface is known.}

Another series of researches corresponds to the classification of the systems of curves of the surface made according to the \emph{genus} $\pi$. For the first values of $\pi$ there are results which it is well to recall.

1) A well known theorem of M.\ \emph{Nöther} relates to the case $\pi = 0$${}^{*)}$:

\emph{Every surface containing a linear system $\infty^{1}$ of rational curves $(\pi = 0)$ is rational.}

In particular: \emph{every surface of ordinary space whose plane sections are rational is rational}; and one may add (thanks to a theorem of M.\ \emph{Picard}${}^{**)}$) that \emph{the surface is ruled, unless it be the surface of Steiner} (of the fourth order, with a triple point and three double straight lines).

2) If $\pi = 1$, the knowledge of a system $\infty^{1}$ of curves does not suffice to conclude from it the rationality of the surface (one may verify this on numerous examples); but if one knows a linear system $\infty^{2}$ of elliptic curves $(\pi = 1)$, one may state an interesting result${}^{***)}$:

\emph{Every surface containing a linear system $\infty^{2}$ of elliptic curves may be transformed birationally into a plane, or else into an elliptic ruled surface.} In this last case there exists on the given surface an elliptic pencil of rational curves which cut each curve of the system $\infty^{2}$ in a single variable point; it is the pencil corresponding to the series of the generators of the ruled surface.

\textbf{*)} l.\ c.\ Math.\ Annalen, 3.

\textbf{**)} \emph{Sur les surfaces algébriques dont toutes les sections planes sont unicursales}, Journal für Mathem., 100, 1885; see also \emph{Guccia}, Rendic.\ Circolo Mat.\ di Palermo, t.\ I, 1886.

\textbf{***)} \emph{Castelnuovo}, \emph{Sulle superficie algebriche che contengono una rete di curve iperellittiche}, Rendic.\ della R.\ Acc.\ d.\ Lincei, 1894; the proof is reported by M.\ \emph{Enriques} in the Memoir \emph{Sui sistemi lineari di superficie algebriche}..., n${}^{\text{o}}$ 8, Math.\ Annalen, 46.

\origpage{308}

In particular${}^{*)}$: \emph{every surface of ordinary space whose plane sections are elliptic curves is rational, or else it is a ruled surface.} In the first case the order of the surface cannot exceed 9.

3) For $\pi = 2$ one has a result analogous to the preceding, provided one has regard to the restriction $n > 2$${}^{**)}$:

\emph{Every surface containing a linear system of curves of genus 2 whose dimension $r \geqq 2$, and whose degree $n > 2$, may be transformed birationally into a plane, or else into a ruled surface whose genus may assume the values 2 or 1.} In the case of the ruled surface of genus 2, the generators of the latter are cut in a single point by every curve of the given system. If on the contrary the ruled surface is elliptic, its generators are met in more than one point by the curves of the system $\infty^{2}$.

\emph{Every surface of ordinary space whose plane sections have the genus 2 is rational (and contains a pencil of conics), or else it is ruled}${}^{***)}$.

4) The last theorems extend at once to every surface which contains a system $\infty^{2}$ of \emph{hyperelliptic} curves whose genus $\pi$ is as high as one wishes${}^{\dagger)}$:

\emph{Every surface containing a linear system $\infty^{2}$ of hyperelliptic curves of genus $\pi$, whose characteristic series is non-special, may be transformed birationally into a plane, or else into a ruled surface whose genus has one of the two values $\pi$, 1.} On the subject of this last case there is occasion to make a remark analogous to the preceding.

\emph{Every surface of ordinary space whose plane sections are hyperelliptic curves is rational (and contains a pencil of conics), or else it is ruled.}

5) For $\pi = 3$, the hyperelliptic case excepted, one has up to the present only incomplete results${}^{\dagger\dagger)}$.

\textbf{*)} \emph{Castelnuovo}, \emph{Sulle superficie algebriche le cui sezioni piane sono curve ellittiche}, Rendic.\ della R.\ Acc.\ d.\ Lincei, 1894; see also \emph{Enriques}, \emph{Sui sistemi lineari}..., n${}^{\text{o}}$ 5.

\textbf{**)} \emph{Castelnuovo}, \emph{Sulle superficie algebriche che contengono una rete}..., l.\ c.; \emph{Enriques} l.\ c.\ n${}^{\text{o}}$ 8.

\textbf{***)} \emph{Enriques}, \emph{Sui sistemi lineari di superficie algebriche}..., Rendic.\ della R.\ Acc.\ d.\ Lincei 1893, Math.\ Annalen 39.

\textbf{†)} See for these results the cited Memoirs of MM.\ \emph{Castelnuovo} and \emph{Enriques}.

\textbf{††)} \emph{Castelnuovo}, \emph{Sulle superficie algebriche le cui sezioni sono curve di genere 3}, Atti dell'Acc.\ d.\ Scienze di Torino, vol.\ 25, 1890.

\origpage{309}

\subsection*{40. Some results which are brought back to the preceding ones.}

If one knows on a surface a sufficiently wide linear system of curves, one may sometimes deduce from it another which is in the conditions considered above, by the aid of a method which has been employed with success by M.\ \emph{Del Pezzo}${}^{*)}$, and which consists in imposing new base points (simple, double...) on the curves of the system. From the projective point of view this amounts to considering the surface of the space $S_{r}$ (of $r$ dimensions) whose sections by the $S_{r-1}$ form the given system, and to projecting the surface from some of its points, of its tangent planes, etc.

By this means one sees at once that \emph{every surface of the space $S_{r}(r \geqq 3)$ having the order $n = r - 1$, or $n = r$, is rational, except the elliptic cone which may present itself in the second case;} one may even assign interesting projective properties of these surfaces${}^{**)}$.

Another result, obtained by the same method${}^{***)}$, completes in a certain sort the theorem cited above, according to which a linear system of genus $\pi > 1$ on a rational surface has the dimension $\leqq 3\pi + 5$:

\emph{If a linear system of curves of genus $\pi \geqq 1$ on a surface has the dimension $r \geqq 3\pi + 5$, the surface is rational, or else it may be represented on a ruled surface of genus $\pi$ whose generators meet the curves of the system in a single point. It is the second case that always presents itself if $r > 3\pi + 5$ for $\pi > 1$, or $r > 9$ for $\pi = 1$.}

One might even, as it appears, obtain more expressive results in relation with lower values of $r$.

\subsection*{41. Method of the successive adjoint systems.}

The method which we have just set out does not lend itself, in most cases, to examining whether a surface possessing a given system of curves is rational. The imposition of new base points on the system will ordinarily lead to a new system whose dimension will result as small as one wishes, but whose degree and genus

\textbf{*)} \emph{Sulle superficie di ordine $n$ immerse nello spazio di $n + 1$ dimensioni} (Rendic.\ Accad.\ di Napoli, 1885); \emph{Sulle superficie dell'n${}^{\text{o}}$ ordine immerse nello spazio di $n$ dimensioni} (Rendic.\ Circolo Mat.\ di Palermo t.\ I, 1887).

\textbf{**)} \emph{Del Pezzo}, l.\ c.

\textbf{***)} \emph{Enriques}, \emph{Sulla massima dimensione}..., Atti dell'Acc.\ d.\ Scienze di Torino, vol.\ 29.

\origpage{310}
will still be too great for one to be able to draw any conclusion from them. Now a procedure of \emph{reduction} cannot satisfy us if it does not necessarily lead to a system of sufficiently small genus or degree every time that the surface considered is rational; for only in that case will it permit us to decide whether a given surface is rational, or else is not. Such a procedure is furnished us by the operation of \emph{adjunction}, which we defined in a preceding Chapter.

Let $|C|$ be a linear system of curves on the given surface; let us form the adjoint system $|C'|$, then the system $|C''|$ adjoint to the latter, which we call the \emph{second adjoint} of $|C|$, and so on. Now, with respect to the succession of the systems $|C|$, $|C'|$, $|C''|$, $\ldots$, two cases may present themselves.

1) Either the succession has an end; one arrives therefore at a last system $|C^{(i)}|$ which admits no further adjoint system;

2) or else the succession is unlimited; each of the systems which compose it admits the adjoint system.

It is naturally the first case that will present itself when the surface considered is rational; one sees it at once if one operates on systems of plane curves. But one cannot say inversely that the condition 1) is always \emph{sufficient} to affirm the rationality of the surface; for one must further add the condition of \emph{regularity} of the surface $(p_{g} = p_{n})$. In fact the last system $|C^{(i)}|$ will then have (as is proved) the genus 0 or 1, and one will be able to apply the results of n${}^{\text{o}}$ 39.

In conclusion, to recognize the rational character of a surface one must first verify that it is regular (for this it suffices in general to examine whether the series cut out by $|C'|$ on the curve $C$ is complete (n${}^{\text{o}}$ 28)).

Next one must verify whether the succession $|C|$, $|C'|$, $|C''|$ $\ldots$ has an end. This procedure requires, as it appears, a number of operations which one cannot assign \emph{a priori}, before leading to a definitive answer on the nature of the surface proposed. This is a serious defect, which we shall try to avoid. One succeeds in it by remarking that all the adjoint systems may be deduced from the two first $|C|$ and $|C'|$ by simple operations of addition and of subtraction; one has in fact (in virtue of the fundamental theorem n${}^{\text{o}}$ 16)
\[ |C''| = |2C' - C|, \quad |C'''| = |3C' - 2C|, \quad \text{etc.} \]
It is therefore the consideration of the systems $|C|$ and $|C'|$ that suffices to settle the question of the rationality of the surface. Here is how one will have to proceed.

a) One must first recognize that $|C'|$ does not contain $|C|$;

\origpage{311}

b) it is further necessary that $|2C'|$ should not contain $|2C|$; this condition includes the preceding, but does not flow from it.

Every ulterior verification is superfluous; in fact it is proved that if the last condition is fulfilled the succession $|C|$, $|C'|$, $|C''|$ $\ldots$ has an end; consequently the surface is rational. One has therefore the following theorem${}^{*)}$:

\emph{On a regular surface a linear system $|C|$ of curves is supposed known} (at least $\infty^{3}$), \emph{as well as its adjoint system $|C'|$; in order that the surface should be rational, it is necessary and sufficient that the system $|2C'|$ should not contain the system $|2C|$.}

We shall see at once some corollaries of this theorem; but let us remark first that the condition named is certainly fulfilled when \emph{the degree $n$ of the system $|C|$ is greater than $2\pi - 2$ ($\pi$ being the genus of $|C|$).}

\subsection*{42. Rationality of a surface deduced from the knowledge of its numerical invariants.}

The most important consequence that may be drawn from the theorem answers the following question: will the class of the rational surfaces be entirely defined by particular values of the invariant numbers $(p_{g}, p_{n}, P_{2}, P_{3} \ldots)$? \emph{Clebsch} had posed himself the analogous question for the class of the rational curves; he resolved it by remarking that a curve is rational if its genus is null, and \emph{viceversa}. We shall see that for surfaces the answer is analogous, although there is more than one character to be considered.

Let us recall first that for a rational surface one has
\[ p_{n} = p_{g} = P_{2} = P_{3} = \cdots = 0; \]
this is what results from the definitions of these characters. These conditions are not, however, all independent, but neither must it be thought that one of them at once entails the others (as happens for curves). In fact the irrational ruled surfaces, for example, have
\[ p_{g} = P_{2} = P_{3} = \cdots = 0, \]
but $p_{n} < 0$; and one likewise knows non-rational surfaces with $p_{n} = 0$. One was really led to believe that the two conditions $p_{n} = p_{g} = 0$ would suffice to characterize the class of the rational surfaces; and it is only when contrary examples were found that one saw that it was quite otherwise; these are the examples which we have already cited in n${}^{\text{o}}$ 30 $(p_{n} = p_{g} = 0, P_{2} > 0)$.

It is the theorem of the preceding paragraph that furnishes us the

\textbf{*)} \emph{Castelnuovo}, \emph{Sulle superficie di genere zero}.

\origpage{312}
answer to our question; for the hypotheses made in it translate into the relations $p_{g} = p_{n}$, $P_{2} = 0$ (which flow from $p_{n} = 0$, $P_{2} = 0$). One has therefore the result${}^{*)}$:

\emph{In order that a surface should be rational, it is necessary and sufficient that the numerical genus $p_{n}$ and the bigenus $P_{2}$ should vanish} $(p_{n} = P_{2} = 0)$.

\subsection*{43. Determination of the rational surfaces of a given order.}

One may make an application of the last theorem to the determination of all the rational surfaces of the order $n$ in ordinary space; this, at least, in theory, for the practical execution requires notions which one does not possess for the moment. The question is reduced to the analysis of the singularities (multiple curves and points) which may be attributed to a surface of the order $n$, having regard to the conditions: 1) that the numerical genus of the surface be null, 2) that there exist no bi-adjoint surfaces of the order $2(n-4)$, besides those which are formed by the given surface together with another surface of order $n - 8$. Thus the difficulty is brought back to determining the influence exercised by a known singularity on the numerical genus of a surface, and the behaviour of a bi-adjoint surface at the same singularity. This is what one may do in the simplest cases; thus:

\emph{If a surface $F$ of the order $n$ possesses a double curve having the order $d$ and the genus $\pi$, and possesses moreover $t$ triple points which are triple also for the said curve (and has no other singularities), the conditions of rationality of the surface are the two following:}

1) \emph{one must have the relation}
\[ \frac{(n-1)(n-2)(n-3)}{6} - (n-4)d + 2t + \pi - 1 = 0, \]

2) \emph{there must exist no surfaces of order $2(n-4)$ which pass doubly through the double curve of $F$ (besides the surfaces which contain the $F$).}

\subsection*{44. Surfaces the points of which have coordinates expressible as rational functions of two parameters.}

The fundamental theorem of n${}^{\text{o}}$ 41 permits us also to answer a question which presents itself first of all in the study of the rational surfaces. The definition which we have given of these surfaces translates into the two following conditions:

\textbf{*)} \emph{Castelnuovo}, l.\ c.

\origpage{313}

1) the coordinates $x$, $y$, $z$ of the general point of the surface are expressed as rational functions of two parameters $\alpha$, $\beta$ (which may be regarded as the coordinates of a point of the plane),
\[ x = f_{1}(\alpha, \beta), \quad y = f_{2}(\alpha, \beta), \quad z = f_{3}(\alpha, \beta); \tag{1} \]

2) inversely the parameters $\alpha$, $\beta$ are expressed rationally as functions of $x$, $y$, $z$,
\[ \alpha = \varphi_{1}(x, y, z), \quad \beta = \varphi_{2}(x, y, z). \tag{2} \]

Now let us suppose that, for a given surface, the first condition is fulfilled, in other terms let us suppose that the (1) subsist; and let us try to deduce the (2) from the (1) by operations of elimination. If one can effect the passage, which happens in many cases, the surface is rational, according to the definition. But it may happen, on the contrary, that one cannot deduce the (2) from the (1), because to a point $(x, y, z)$ of the surface there correspond $n > 1$ points $(\alpha, \beta)$ of the plane. Will it be necessary to conclude in this case that the surface is not rational? Or else, on the contrary, will one be able to introduce, in place of $\alpha$ and $\beta$, two new parameters $\alpha'$, $\beta'$ (rational functions of the first), in such a way that $x$, $y$, $z$ are expressed rationally by the aid of $\alpha'$, $\beta'$, and \emph{viceversa} $\alpha'$, $\beta'$ are expressed rationally by the aid of $x$, $y$, $z$? In this last case the surface would naturally be rational, and the peculiarity just remarked with regard to $\alpha$ and $\beta$ would depend only on the special choice of the parameters.

Now it was indeed suspected that it was so; for in the same sense the analogous question relating to curves had been resolved by M.\ \emph{Lüroth}. But the proof of M.\ \emph{Lüroth}, as it appears, cannot be extended to surfaces.

Here is how one succeeds in proving the theorem.

To every point $(x, y, z)$ of the surface there correspond (we have said) $n > 1$ points $(\alpha, \beta)$ of the plane; to every algebraic curve of the surface there corresponds an algebraic curve of the plane, and to a linear system of curves on the former, a linear system of curves on the latter. Now from the known properties of these last systems one may deduce properties of the first systems. Thus one sees for example that on the surface every complete linear system has its characteristic series complete, whence it results that the two genera $p_{g}$, $p_{n}$ of the surface are equal. One sees, moreover, that on the surface there exist systems of genus $\pi$ with degree $> 2\pi - 2$, whence it follows
\[ p_{g} = P_{2} = P_{3} = \cdots = 0. \]
The two conditions $(p_{n} = P_{2} = 0)$ for the rationality of a surface are therefore fulfilled. This is what the following theorem affirms${}^{*)}$:

\textbf{*)} \emph{Castelnuovo}, \emph{Sulla razionalità delle involuzioni piane}, Math.\ Ann.\ 44, 1893.

\origpage{314}

\emph{If the coordinates of the general point of a surface can be expressed rationally as functions of two parameters, the surface is rational.}

One may present this theorem under another form, provided one remarks that the groups of $n$ points $(\alpha, \beta)$ corresponding to the points $(x, y, z)$ of our surface form on the plane an \emph{involution}, that is to say a system $\infty^{2}$ of groups, each group of which is entirely determined by one of its points:

\emph{Every involution on the plane is rational.}

\subsection*{45. Rationality of a surface deduced from the knowledge of a series, not linear, of curves.}

Up to now we have confined ourselves to considering linear systems of curves. Now there are cases where the knowledge of a non-linear series $\infty^{1}$ of curves on a surface, a series such that through the general point of the surface there passes more than one curve of the series, may suffice to recognize the rationality of the latter. Here are some examples:

a) A first example regards the case where the curves of the series are rational, as well as the series itself (considering in it the curve as element): one has then the following theorem, which is deduced at once from the theorem on the plane involutions (just as the latter follows from the former):

\emph{Every surface which contains a rational series of rational curves is rational}${}^{*)}$.

This theorem comprises, as a particular case, the theorem of M.\ \emph{Nöther} relative to the surfaces having a \emph{linear} system $\infty^{1}$ of rational curves; to recover it, it suffices in fact to suppose that through every point of the surface there passes a single curve of the series.

b) A second theorem considers a series whose degree (the number of variable intersections of two curves) is 1:

\emph{Every surface which contains an algebraic series of curves whose degree is 1, and such that through each point of the surface there pass $i > 2$ curves of the series, is rational}${}^{**)}$.

If, the other hypotheses remaining the same, one supposes $i = 2$, the surface may be represented birationally on the variety formed

\textbf{*)} \emph{Castelnuovo}, \emph{Sulle superficie algebriche che contengono una rete di curve iperellitiche}, l.\ c.

\textbf{**)} \emph{Humbert}, Comptes Rendus, 1893; \emph{Sur quelques points de la théorie}..., Journal de Mathématiques 4${}^{\text{e}}$ s${}^{\text{e}}$, t${}^{\text{e}}$ X; \emph{Castelnuovo}, \emph{Sulla linearità delle involuzioni}..., Atti dell'Acc.\ d.\ Sc.\ di Torino, vol.\ 28, 1893.

\origpage{315}
by the couples of points of a curve whose genus is equal to that of the series.

c) To the preceding theorem is attached the following${}^{*)}$:

\emph{Let there exist on a surface two algebraic series of curves such that two general curves belonging to different series meet in a single variable point; then if through the general point of the surface there passes more than one curve of each of the series, the surface is rational, and the curves of the series are themselves rational.} If through the general point of the surface there passes a single curve of the first series, and more than one of the second, the surface may be transformed birationally into a ruled surface; if finally through the point there passes a single curve of each series, the surface may be represented on the variety of the couples of points of two curves.

\subsection*{46. Surfaces admitting a continuous group of birational transformations into themselves.}

To complete, as far as possible, this review of the results obtained on the subject of the rational surfaces, we must say a few more words on the surfaces which admit a continuous group of birational transformations into themselves, from the point of view of their rational character.

M.\ \emph{Picard}${}^{**)}$ was the first to approach the problem of the groups of transformations of a surface into itself, and he arrived at extremely remarkable results which do not find their place here, for they concern principally certain classes of irrational surfaces. For our purpose it suffices in fact to state the following theorems, which give the natural extension of the well known result of M.\ \emph{Schwarz} relative to curves:

\emph{Every surface which admits a continuous transitive group of birational transformations into itself, a group depending on two parameters, is rational, if the transformations are not interchangeable two by two.}

\emph{Every surface which admits a continuous transitive group of birational transformations into itself, a group depending on more than two parameters, is rational, or else may be transformed into an elliptic ruled surface}${}^{***)}$.

\textbf{*)} \emph{Castelnuovo}, l.\ c.

\textbf{**)} See the Memoirs cited in the Introduction.

\textbf{***)} \emph{Castelnuovo} and \emph{Enriques}, Comptes Rendus de l'Ac.\ d.\ Sc.\ 1895; see also in the same collection (1895) a Note of M.\ \emph{Painlevé}. On the subject of the statements which we report here, it is well to remark that we call \emph{transitive} (after M.\ Lie) a group which can change a point into another point arbitrarily fixed; and that the second theorem says a little more than the one which is found in the Comptes Rendus.

\origpage{316}

\subsection*{47. Double planes of genus one.}

In n${}^{\text{o}}$ 38 we have spoken of the \emph{double planes}, that is to say of the algebraic fields
\[ x, y, \sqrt{f(x, y)}, \]
where $f$ designates a polynomial; we have stated the results of \emph{Clebsch} and \emph{Nöther} on the nature of the branch curve $f(x, y) = 0$ when the double plane is rational, and in consequence $\sqrt{f(x, y)}$ becomes a rational function of two new parameters, by a suitable rational substitution on $x$, $y$.

The problem now presents itself of classifying the higher double planes (according to their branch curves). The results obtained up to now on this subject are confined to the case of the double planes whose genera $p_{n}$ and $P_{2}$ (and consequently $p^{(1)}$, $p_{g}$, $P_{3} \ldots$) have the value 1. One has the theorem${}^{*)}$:

\emph{The double planes whose genera $(p_{n}, P_{2})$ have the value 1 are distributed into four classes, according as their branch curve can be brought back, by a birational transformation of the plane, to one of the following curves:}

1) \emph{a general curve of the order 6:}

2) \emph{a curve of the order 8 endowed with two points of the fourth order;}

3) \emph{a curve of the order 10 endowed with a point of the order 7 and with two triple points infinitely near to it;}

4) \emph{a curve of the order 12 endowed with a point of the order 9 and with three triple points infinitely near to it.}

\emph{Florence}, 13 February 1896.

\textbf{*)} \emph{Enriques}, \emph{Sui piani doppi di genere uno}, l.\ c.

\end{document}
